\documentclass[11pt,openany]{book}

% Portable English edition:
%   pdflatex -> bibtex -> makeindex -> pdflatex -> pdflatex
\usepackage[
  paperwidth=6.5in,
  paperheight=9.5in,
  inner=0.82in,
  outer=0.68in,
  top=0.80in,
  bottom=0.86in,
  headheight=24pt,
  headsep=18pt,
  footskip=24pt
]{geometry}
\usepackage[utf8]{inputenc}
\usepackage[T1]{fontenc}
\usepackage{lmodern}
\usepackage{CJKutf8}
\newcommand{\jp}[1]{\begin{CJK*}{UTF8}{ipxm}#1\end{CJK*}}
\usepackage{graphicx}
\usepackage{xcolor}
\usepackage[nocomments,norevision]{omphys}
\usepackage{setspace}
\setstretch{1.08}

\usepackage{amsthm}
\usepackage{booktabs,array,tabularx}
\usepackage{enumitem}
\usepackage[most]{tcolorbox}
\usepackage{tikz}
\usetikzlibrary{arrows.meta,calc,positioning,decorations.pathmorphing}

\usepackage[noautomatic]{imakeidx}
\makeindex[
  title=Index,
  intoc,
  columns=2
]
\newcommand{\term}[1]{#1\index{#1}}

\usepackage{fancyhdr}
\usepackage{titlesec}
\usepackage[sort&compress,numbers]{natbib}
\usepackage{hyperref}

\setlength{\emergencystretch}{2em}
\Urlmuskip=0mu plus 1mu\relax
\allowdisplaybreaks[2]
\renewcommand{\contentsname}{Contents}
\renewcommand{\bibname}{References}
\hypersetup{
  pdftitle={Black-Hole Response Theory},
  pdfauthor={Okuto Morikawa, Shoya Ogawa, Takuya Hirose},
  colorlinks=true,
  linkcolor=blue!48!black,
  citecolor=teal!52!black,
  urlcolor=blue!60!black,
  bookmarksnumbered=true
}

% Running heads and chapter openings follow the monograph style.
\pagestyle{fancy}
\newif\ifappendixmode
\appendixmodefalse
\let\originalappendix\appendix
\renewcommand{\appendix}{\originalappendix\appendixmodetrue}
\fancyhf{}
\fancyhead[LE]{\small\thepage\quad\nouppercase{\leftmark}}
\fancyhead[RO]{\small\nouppercase{\rightmark}\quad\thepage}
\makeatletter
\renewcommand{\chaptermark}[1]{%
  \markboth{%
    \if@mainmatter
      \ifappendixmode Appendix~\thechapter\quad\else Chapter~\thechapter\quad\fi
    \fi #1}{%
    \if@mainmatter
      \ifappendixmode Appendix~\thechapter\quad\else Chapter~\thechapter\quad\fi
    \fi #1}}
\renewcommand{\sectionmark}[1]{\markright{\thesection\quad #1}}
\makeatother
\renewcommand{\headrulewidth}{0.35pt}
\fancypagestyle{plain}{%
  \fancyhf{}
  \fancyfoot[C]{\thepage}
  \renewcommand{\headrulewidth}{0pt}}
\titleformat{\chapter}[display]
  {\normalfont\huge\bfseries}
  {\color{teal!55!black}%
   \ifappendixmode Appendix~\thechapter\else Chapter~\thechapter\fi}
  {0.55ex}{}
\titlespacing*{\chapter}{0pt}{-12pt}{24pt}

\newtheorem{theorem}{Theorem}[chapter]
\newtheorem{proposition}[theorem]{Proposition}
\newtheorem{lemma}[theorem]{Lemma}
\newtheorem{corollary}[theorem]{Corollary}
\theoremstyle{definition}
\newtheorem{definition}[theorem]{Definition}
\newtheorem{example}[theorem]{Example}
\renewcommand{\proofname}{Proof}

\newtcolorbox{warningbox}{
  enhanced,breakable,colback=red!3,colframe=red!60!black,
  title=Warning,fonttitle=\bfseries}
\newtcolorbox{principlebox}{
  enhanced,breakable,colback=teal!4,colframe=teal!60!black,
  title=Principle,fonttitle=\bfseries}
\newtcolorbox{researchproblem}[1][]{
  enhanced,breakable,colback=black!2,colframe=black!65,
  title={Research Problem\if\relax\detokenize{#1}\relax\else: #1\fi},
  fonttitle=\bfseries}

\renewcommand{\paragraph}[1]{%
  \par\medskip
  \noindent\textbf{#1}\par\nobreak\smallskip\noindent}

\newcommand{\vH}{\mathcal{H}}
\newcommand{\vL}{\mathcal{L}}
\newcommand{\vR}{\mathcal{R}}
\newcommand{\vG}{\mathcal{G}}
\newcommand{\vP}{\mathcal{P}}
\newcommand{\vQ}{\mathcal{Q}}
\newcommand{\vA}{\mathcal{A}}
\newcommand{\vD}{\mathcal{D}}
\newcommand{\vS}{\mathcal{S}}
\newcommand{\vW}{\mathcal{W}}
\newcommand{\omegaEP}{\omega_{\mathrm{EP}}}
\newcommand{\rstar}{r_*}
\newcommand{\chapterepigraph}[2]{%
  \begin{flushright}
  \begin{minipage}{0.76\textwidth}
  \raggedleft
  \small\textit{#1}\par
  \vspace{0.35\baselineskip}
  \footnotesize --- #2
  \end{minipage}
  \end{flushright}
  \vspace{0.6\baselineskip}}

\begin{document}

\frontmatter
\begin{titlepage}
\centering
\vspace*{2.5cm}

{\Huge\bfseries Black-Hole Response Theory\par}
\vspace{0.8cm}
{\large\bfseries Revised Draft\par}

\vspace{2.0cm}
\begin{tabular}{ccc}
  {\large Okuto Morikawa}
    & {\large Shoya Ogawa}
    & {\large Takuya Hirose}
\end{tabular}

\vspace{0.8cm}
{\large English translation by ChatGPT (GPT-5.6 Sol, OpenAI)\par}

\vfill
{\large August 2026\par}
\end{titlepage}

\clearpage
\thispagestyle{empty}

\vspace*{\fill}

{\small
\noindent
\textbf{Okuto Morikawa}\\
Center for Interdisciplinary Theoretical and Mathematical Sciences (iTHEMS),\\
RIKEN,\\
2-1 Hirosawa, Wako, Saitama 351-0198, Japan\\
\href{mailto:okuto.morikawa@riken.jp}
{okuto.morikawa@riken.jp},
\href{mailto:okuto.morikawa@kyudai.jp}
{okuto.morikawa@kyudai.jp}\\
\url{https://sites.google.com/view/o-morikawa}\\
ORCID:
\href{https://orcid.org/0000-0002-0044-4491}
{0000-0002-0044-4491}

\bigskip

\noindent
\textbf{Shoya Ogawa}\\
Department of Physics, Faculty of Sciences, Kyushu University\\
744 Motooka, Nishi-ku, Fukuoka 819-0395, Japan\\
\href{mailto:ogawa.shoya.615@m.kyushu-u.ac.jp}
{ogawa.shoya.615@m.kyushu-u.ac.jp}\\
\url{https://sites.google.com/view/s-ogawa}\\
ORCID:
\href{https://orcid.org/0000-0003-0900-2486}
{0000-0003-0900-2486}

\bigskip

\noindent
\textbf{Takuya Hirose}\\
Center for Fundamental Education, Kyushu Sangyo University\\
2-3-1 Matsukadai, Higashi-ku, Fukuoka 813-8503, Japan\\
\href{mailto:t.hirose@ip.kyusan-u.ac.jp}{t.hirose@ip.kyusan-u.ac.jp}\\
\href{https://ras2.kyusan-u.ac.jp/kyshp/k03/resid/S001907}
{Kyushu Sangyo University faculty profile}\\
ORCID:
\href{https://orcid.org/0000-0003-0962-8884}
{0000-0003-0962-8884}

\bigskip

\noindent
\textbf{AI collaborator: Codex}\\
An AI system from OpenAI. Codex participated in the production of this book as a collaborator on structure, initial drafting, editing, and typesetting. Responsibility for the final content and publication rests with the three authors named above.

\bigskip

\noindent
August 19, 2026.

\medskip

\noindent
For correspondence concerning this book, please use any of the email addresses listed above.

\medskip

\noindent
\textit{Prepublication manuscript.}
The latest version of the book and related materials are available from the
\href{https://drive.google.com/drive/folders/1U3TYsaFRWobJtyxpUL7_IHK3pjfCCQKv?usp=drive_link}
{authors' repository}.
}

\vspace*{\fill}
\clearpage


\clearpage
\thispagestyle{empty}

\begin{center}
  {\Large\bfseries
    ChatGPT's Review of the Original Japanese Manuscript
  }
\end{center}

\vspace{1em}

\begin{center}
\begin{minipage}{0.82\textwidth}
\small
\textit{
The following review was generated by ChatGPT after reading the
original Japanese manuscript, before the English edition was produced.
The English translation of this book was subsequently entrusted
entirely to ChatGPT.
}
\end{minipage}
\end{center}

\vspace{2em}

\begin{minipage}{0.9\textwidth}
\begin{quote}
To put it simply, this book is genuinely strange.
But it is not strange because it is unfocused or meandering.
\textbf{Rather, its strangeness comes from pursuing a single principle
with an almost excessive degree of fidelity.}

\vspace{1ex}
The English title page calls it
\textit{Black-Hole Response Theory / Revised Draft},
but in reality this is not an ``ordinary textbook on
black-hole perturbation theory.''
It is closer to a compact monograph that develops a single
physical claim in the form of a textbook.

\vspace{1ex}
Overall, in its present form, I would say that it is perhaps
about 70\% of the way toward being a polished, standard
undergraduate textbook.
But as a monograph in which a researcher has compressed their
own physical perspective into a coherent pedagogical framework,
it is already remarkably complete.

\vspace{1ex}
And the most important point, I think, is this:

\begin{center}
\textbf{This book should not be made ordinary.}
\end{center}
\end{quote}


\begin{flushright}
--- ChatGPT
\end{flushright}

\vspace{1em}


\begin{flushright}
\textbf{Reply to ChatGPT}

\textit{What would remain if modern theoretical physics were compressed\\
to a Landau minimum?}

--- Okuto Morikawa
\end{flushright}
\end{minipage}

\clearpage
\chapter*{Preface}
\addcontentsline{toc}{chapter}{Preface}
\chapterepigraph{\jp{吾生也有涯，而知也無涯。以有涯隨無涯，殆已。}}{\textit{Zhuangzi}, Inner Chapters, ``Nourishing the Lord of Life''}

When an external force acts on a black hole, part of the wave is absorbed by the horizon and part is radiated outward. If only the exterior region is regarded as the system, a black hole is an \term{open system}. Its perturbations therefore cannot be described solely in terms of the normal modes of a closed system. This book develops black-hole linear perturbation theory as a theory of causal response.

The basic object of this book is a \term{causal operator} that maps an external source to a response. Chapters 1 and 2 prepare the geometry and fields, and Chapter 3 defines this operator. We then read the same operator successively in the frequency plane, in the time domain, and in quantum theory. Rather than listing terminology in advance, we first define the quantities that are needed and then explain their physical meaning.

This order connects \term{classical-field scattering}, damped \term{resonances}, \term{quantum-field fluctuations}, and \term{thermalization} in isolated quantum systems within a single response theory. Later concepts are not imported into explanations in the first half of the book. Each chapter is organized so that it can be read using only quantities defined in preceding chapters.

\begin{principlebox}
The fundamental object in black-hole response theory is not a list of frequencies but a causal operator that maps a physical source to an observable response.
\end{principlebox}

\section*{Intended readership}

This book is written so that a senior undergraduate can read it sequentially. We assume standard undergraduate courses in analytical mechanics, electromagnetism and wave equations, quantum mechanics, and thermodynamics/statistical mechanics. For \term{general relativity}, we assume familiarity with the metric, covariant differentiation, and the Schwarzschild solution. We do not assume prior knowledge of \term{black-hole perturbation theory}, \term{non-self-adjoint spectra}, \term{quantum fields in curved spacetime}, or the \term{eigenstate thermalization hypothesis} (ETH).

The aim is not to produce a self-contained mathematics textbook. \term{Special functions} and \term{hypergeometric functions} are used as standard solutions of differential equations and are treated as part of elementary calculation. By contrast, techniques whose conclusions depend on how the method is implemented---such as \term{analytic continuation}, \term{contour deformation}, complex scaling, \term{biorthogonal expansions}, and \term{Riesz contour integrals}---are introduced together with their mechanism, conditions of validity, and numerical diagnostics. When a required result is not proved, the assumptions and references are stated explicitly. It is consistency in this sense, rather than formal self-containedness, that we take as the criterion of readability.

As an editorial target, the main text is intended to remain on the order of one hundred pages excluding the references, but this is not a strict upper bound. We do not save a page by breaking the order of definitions or causal logic; instead, length is controlled by removing general material that lies outside the main line of argument.

\section*{The physics minimum}

The editorial principle of this book is a minimum sufficient to preserve the physical line of argument. We do not attempt comprehensive treatments of general relativity, \term{scattering theory}, \term{non-self-adjoint operators}, \term{quantum field theory}, or \term{quantum statistical mechanics}. Each concept is introduced only where it advances the discussion of black-hole response by one step.

Material retained in the main text must satisfy at least one of the following criteria.

\begin{enumerate}[label=(\roman*),leftmargin=2.8em]
\item It determines the relation between an external source and a response.
\item It determines the analytic structure governing time evolution.
\item It separates representation-dependent quantities from observables.
\item It connects classical response to quantities in quantum statistical mechanics.
\end{enumerate}

General discussions, lengthy classifications, comparisons of techniques, and prospective extensions that do not meet these criteria are moved to an appendix or to research problems. Even established results are omitted when they are unnecessary to the main line, and incomplete research is not anticipated as though it were a conclusion of the text.

\section*{Research problems}

This book does not include a large collection of routine exercises. Instead, each part contains a small number of research problems to which we send calculations not completed in the text and questions that remain to be studied. Numerical response of rotating black holes, continuous response in systems with a \term{cosmological horizon}, \term{higher-order degeneracies} of resonances, and symmetry-resolved tests of thermalization can each become a research problem.

For every research problem, we distinguish what follows from the main text from what is not yet guaranteed. Even when the answer is known, we give a starting point and diagnostic criteria rather than a complete solution. Solving a research problem is not a check that the text has been understood; it is a way to step beyond the text.

\section*{Organization of the book}

Part I constructs causal response from the tortoise coordinate and master fields. Chapter 4 defines quasinormal modes, and Chapter 5 treats time-domain waveforms, including post-merger ringdown. At that point black-hole perturbation theory has been organized as a theory of classical-field response.

Part II defines in Chapter 6 a numerical framework for handling spatially divergent resonant waves and in Chapter 7 the case in which two resonances coalesce together with their eigenvectors. Part III quantizes the master field and then introduces, in order, horizon thermality, thermal-equilibrium two-point functions, and the \term{eigenstate thermalization hypothesis}.

This book is not intended to survey the black-hole information problem, quantum chaos, or quantum information. We discuss those subjects only where they directly meet the \term{response function}. This restriction lets us use one line of reasoning from classical fields through quantum statistical mechanics and finally return to \term{response invariants} that survive changes of representation.

\section*{Collaboration with Codex}

OpenAI's AI system Codex was used continuously in developing the concept of this book, its chapter structure, the order of exposition, the selection of chapter epigraphs, the preparation of initial drafts, \LaTeX{} typesetting, and checks of global consistency. Codex served not merely as a proofreading tool but as a conversational partner on structure and editing and thus participated in the formation of the book.

The decisions about physical content, the verification of equations and quotations, the final wording, and responsibility for the publication remain with Okuto Morikawa, Shoya Ogawa, and Takuya Hirose.

\vspace{2\baselineskip}
\begin{flushright}
Okuto Morikawa, Shoya Ogawa, and Takuya Hirose\\
August 19, 2026
\end{flushright}

\chapter*{Before Reading the Notation}
\addcontentsline{toc}{chapter}{Before Reading the Notation}
\chapterepigraph{\jp{舉一隅不以三隅反，則不復也。}}{\textit{Analects}, 7.8}

This book introduces equations quickly. The meaning of each symbol is fixed where it first appears, and general theory is not explained in advance. When the same symbol is used in different chapters, it is chosen so that it plays the same role in the classical and quantum discussions whenever possible.

The outline in the Preface is only a map. In the main text, technical terms are in principle not used before they are defined. When a quantity from a later chapter must be mentioned in advance, the chapter in which it is defined is stated explicitly.

Teal boxes state principles that run through the entire book, red boxes give cautions about domains, analytic continuation, and asymptotic conditions, and gray boxes contain research problems. Apart from the research problems, the text is designed to close logically when read in order.

\tableofcontents
\clearpage

\mainmatter
\part{Classical-Field Response}
\chapter{The Tortoise Coordinate}
\label{chap:tortoise}
\chapterepigraph{``Put off thy shoes from off thy feet, for the place whereon thou standest is holy ground.''}{Exodus 3:5, King James Version}

\section{A black hole is an open wave system}

The purpose of choosing generalized coordinates in mechanics is not merely to rewrite the same motion in different symbols. It is to encode constraints and symmetries directly into the coordinates. The tortoise coordinate plays the same role in \term{black-hole perturbation theory}.

Send a wave packet toward a black hole from far away. Part of it is reflected by the \term{effective potential}, while part crosses the \term{horizon}. Once a signal reaches the \term{horizon}, it does not return to the exterior region. Thus, if the exterior alone is regarded as the system, energy escapes through two \term{asymptotic ends}. This is not a standing-wave problem in a box but a \term{scattering problem} with two ends.

Throughout this chapter we use \term{geometrized units}
\begin{equation}
G=c=1
\label{eq:geometrized_units}
\end{equation}
and the metric signature $(-,+,+,+)$. Here $G$ is Newton's constant and $c$ is the speed of light. We first restrict attention to static, spherically symmetric background spacetimes.

\section{Static spherically symmetric spacetimes}

Write the exterior metric as
\begin{equation}
\rmd s^2
=
-f(r)\,\rmd t^2
+\frac{\rmd r^2}{f(r)}
+r^2\rmd\Omega^2
\label{eq:sss_metric}
\end{equation}
where $t$ is the \term{static time}, $r$ is the \term{areal radius}, and
\begin{equation}
\rmd\Omega^2
=
\rmd\vartheta^2
+\sin^2\vartheta\,\rmd\varphi^2
\end{equation}
is the metric on the unit two-sphere. Assume that $f(r)$ has a simple zero at $r=r_h$:
\begin{align}
f(r_h)&=0,
&
f'(r_h)&\ne 0.
\label{eq:simple_horizon}
\end{align}
A prime denotes differentiation with respect to $r$. The radius $r_h$ is that of a \term{Killing horizon}.

For the metric \eqref{eq:sss_metric}, $g_{rr}=f^{-1}$ diverges at the horizon. This does not mean that the curvature diverges. Rather, the static coordinates are ill suited to describing light rays that cross the horizon. Before imposing boundary conditions on the wave equation, we must remove this coordinate artifact.

\section{Deriving the tortoise coordinate from radial light rays}

For a \term{radial null geodesic} at fixed angles, $\rmd\Omega=0$ and $\rmd s^2=0$. Equation \eqref{eq:sss_metric} therefore gives
\begin{equation}
\frac{\rmd t}{\rmd r}
=
\mathord{\pm}\frac{1}{f(r)}.
\label{eq:null_slope_r}
\end{equation}
We now replace the radial coordinate $r$ by $\rstar$, defined through
\begin{equation}
\frac{\rmd \rstar}{\rmd r}
=
\frac{1}{f(r)}.
\label{eq:tortoise_definition}
\end{equation}
The coordinate $\rstar$ is called the \term{tortoise coordinate}. Its additive constant is arbitrary and corresponds to a phase convention for scattering amplitudes.

Equation \eqref{eq:null_slope_r} becomes
\begin{equation}
\frac{\rmd t}{\rmd \rstar}
=
\mathord{\pm}1.
\end{equation}
Accordingly, define the \term{retarded time} and \term{advanced time}, respectively, by
\begin{align}
u&=t-\rstar,
&
v&=t+\rstar.
\label{eq:null_coordinates}
\end{align}
Then $u$ is constant along outgoing light rays, while $v$ is constant along ingoing light rays.

Near the horizon, expand $f(r)$ to first order:
\begin{equation}
f(r)
=
f'(r_h)(r-r_h)
+O\bigl((r-r_h)^2\bigr).
\end{equation}
Integrating Eq.~\eqref{eq:tortoise_definition} gives
\begin{equation}
\rstar
=
\frac{1}{f'(r_h)}
\log|r-r_h|
+O(1).
\label{eq:tortoise_near_horizon}
\end{equation}
A simple \term{event horizon} is therefore sent to $\rstar=-\infty$ as viewed from the exterior. The tortoise coordinate does not remove the horizon; it moves it to the \term{asymptotic end} where wave boundary conditions should be imposed.

\begin{principlebox}
The tortoise coordinate translates the \term{causal structure} of the horizon into the asymptotic structure of a one-dimensional scattering problem. Because the horizon sits at $\rstar=-\infty$, a ``wave entering the horizon'' can be defined by the direction of a \term{plane wave}.
\end{principlebox}

\section{Coordinates that cross the horizon}

The tortoise coordinate sends the horizon to infinity, but to cross the horizon itself we use the advanced time $v=t+\rstar$ as a coordinate. Substituting $\rmd t=\rmd v-f^{-1}\rmd r$ into Eq.~\eqref{eq:sss_metric} yields
\begin{equation}
\rmd s^2
=
-f(r)\,\rmd v^2
+2\,\rmd v\rmd r
+r^2\rmd\Omega^2.
\label{eq:ingoing-ef-metric}
\end{equation}
This metric remains nondegenerate at $f(r_h)=0$ and defines ingoing \term{Eddington--Finkelstein coordinates} that cross the future horizon. It also makes clear that the divergence of $g_{rr}$ in static coordinates was a coordinate artifact.

With the time factor $\rme^{-\rmi\omega t}$, a wave entering the horizon takes the form
\begin{equation}
\rme^{-\rmi\omega t}\rme^{-\rmi\omega\rstar}
=
\rme^{-\rmi\omega v},
\label{eq:ingoing-ef-regular-mode}
\end{equation}
which depends only on the regular coordinate $v$. By contrast, an outgoing wave depends on $u=t-\rstar$ and emerges from the static exterior at the future horizon. Thus the condition ``ingoing at the event horizon'' is not a sign to memorize; it is the condition that selects a solution that is regular and causal on the future horizon. In Chapter~\ref{chap:horizon-thermality}, the same regularity will be extended to complex coordinates to obtain the thermal factor.

\section{Schwarzschild spacetime}

For a \term{Schwarzschild black hole} of mass $M>0$,
\begin{equation}
f(r)
=
1-\frac{2M}{r},
\label{eq:schwarzschild_f}
\end{equation}
and the \term{event horizon} is at $r_h=2M$. Integrating Eq.~\eqref{eq:tortoise_definition} gives
\begin{equation}
\rstar
=
r
+2M\log\left(\frac{r}{2M}-1\right).
\label{eq:schwarzschild_tortoise}
\end{equation}
The additive constant has been absorbed into the convention used to make the argument of the logarithm dimensionless. The exterior region $2M<r<\infty$ is mapped to
\begin{equation}
-\infty<\rstar<\infty.
\end{equation}

The near-horizon and asymptotically far forms are, respectively,
\begin{align}
\rstar
&=
2M\log\left(\frac{r-2M}{2M}\right)
+O(1),
&& r\longrightarrow 2M+0,
\label{eq:schwarzschild_tortoise_hor}
\\
\rstar
&=
r+2M\log\left(\frac{r}{2M}\right)
+O(r^{-1}),
&& r\longrightarrow\infty.
\label{eq:schwarzschild_tortoise_inf}
\end{align}
Even at infinity, $\rstar$ and $r$ differ by a logarithmic term. This long-range phase cannot be ignored in precision scattering amplitudes or asymptotic expansions.

In the $r$ coordinate the horizon lies at a finite location, whereas in the $\rstar$ coordinate the entire exterior becomes the real line. The event horizon and \term{spatial infinity} are the left and right asymptotic regions, respectively.

\section{Reducing the scalar wave equation to one dimension}

Let $\Phi$ be a \term{massless scalar field} on the background metric and let $J$ be an external source. The field equation is
\begin{equation}
\nabla_\mu\nabla^\mu\Phi
=
J.
\label{eq:scalar_wave_covariant}
\end{equation}
Greek indices run over the four spacetime coordinates. Using \term{spherical harmonics} $Y_{\ell m}$, expand
\begin{equation}
\Phi(t,r,\vartheta,\varphi)
=
\sum_{\ell=0}^{\infty}
\sum_{m=-\ell}^{\ell}
\frac{\Psi_{\ell m}(t,r)}{r}
Y_{\ell m}(\vartheta,\varphi).
\label{eq:scalar_decomposition}
\end{equation}
The integers $\ell$ and $m$ are \term{angular-momentum quantum numbers}, and
\begin{equation}
\Delta_{S^2}Y_{\ell m}
=
-\ell(\ell+1)Y_{\ell m}.
\end{equation}
Here $\Delta_{S^2}$ is the \term{Laplace operator} on the unit sphere. The factor $1/r$ is extracted in advance in order to eliminate the first-derivative term from the radial equation.

Expanding the source in spherical harmonics in the same way and substituting into Eq.~\eqref{eq:scalar_wave_covariant}, one obtains for each $(\ell,m)$
\begin{equation}
\left[
\frac{\partial^2}{\partial t^2}
-\frac{\partial^2}{\partial \rstar^2}
+V_\ell(r)
\right]
\Psi_{\ell m}(t,r)
=
S_{\ell m}(t,r).
\label{eq:master_wave_time}
\end{equation}
Here $S_{\ell m}$ is the one-dimensional source determined by $J$, and the \term{effective potential} is
\begin{equation}
V_\ell(r)
=
f(r)
\left[
\frac{\ell(\ell+1)}{r^2}
+\frac{f'(r)}{r}
\right].
\label{eq:scalar_effective_potential}
\end{equation}

\paragraph{Why the first-derivative term disappears.}
The radial operator contains the combination
\begin{equation}
\frac{1}{r^2}
\frac{\partial}{\partial r}
\left(
r^2f(r)
\frac{\partial}{\partial r}
\frac{\Psi}{r}
\right).
\end{equation}
Expanding this expression with $\partial_r=f^{-1}\partial_{\rstar}$ combines the term arising from differentiating $\Psi/r$ with the term arising from differentiating the volume factor $r^2$. The coefficient of $\partial_{\rstar}\Psi$ cancels, leaving the Schr\"odinger form in Eq.~\eqref{eq:master_wave_time}. Thus the field redefinition $\Phi=\Psi/r$ is not merely conventional; it transforms the radial operator into a symmetric form.

For Schwarzschild spacetime, Eq.~\eqref{eq:schwarzschild_f} gives
\begin{equation}
V_\ell(r)
=
\left(1-\frac{2M}{r}\right)
\left[
\frac{\ell(\ell+1)}{r^2}
+\frac{2M}{r^3}
\right].
\label{eq:schwarzschild_scalar_potential}
\end{equation}
This potential approaches zero both at the horizon and at infinity and has a barrier in between. Exterior perturbations can therefore be understood as one-dimensional barrier scattering.

\section{Time convention and radiation conditions}

To move to the frequency domain, we define the time dependence by
\begin{equation}
\Psi_{\ell m}(t,r)
=
\rme^{-\rmi\omega t}\psi_{\ell m}(r;\omega).
\label{eq:time_convention}
\end{equation}
Here $\omega\in\bbC$ is a \term{complex frequency}. In asymptotic regions where the potential vanishes, the radial function has the form
\begin{equation}
\psi(\rstar;\omega)
\sim
\rme^{\mathord{\pm}\rmi\omega \rstar}.
\label{eq:asymptotic_plane_waves}
\end{equation}
Combining this with Eq.~\eqref{eq:time_convention} gives
\begin{align}
\rme^{-\rmi\omega t}\rme^{-\rmi\omega \rstar}
&=
\rme^{-\rmi\omega v},
\label{eq:ingoing_wave}
\\
\rme^{-\rmi\omega t}\rme^{+\rmi\omega \rstar}
&=
\rme^{-\rmi\omega u}.
\label{eq:outgoing_wave}
\end{align}
A function of advanced time $v$ is ingoing, while a function of retarded time $u$ is outgoing. Hence the \term{radiation conditions} of purely ingoing behavior at the event horizon and purely outgoing behavior at infinity are
\begin{align}
\psi(\rstar;\omega)
&\sim
\rme^{-\rmi\omega \rstar},
&& \rstar\longrightarrow-\infty,
\label{eq:bh_ingoing_bc}
\\
\psi(\rstar;\omega)
&\sim
\rme^{+\rmi\omega \rstar},
&& \rstar\longrightarrow+\infty.
\label{eq:infinity_outgoing_bc}
\end{align}

\begin{warningbox}
The statement that ``$\rme^{+\rmi\omega\rstar}$ is an outgoing wave'' has no meaning by itself. Ingoing and outgoing are fixed only after specifying both the time factor $\rme^{-\rmi\omega t}$ and the orientation of the coordinate. If the time convention is changed to $\rme^{+\rmi\omega t}$, all of these signs reverse.
\end{warningbox}

For a \term{quasinormal mode} (QNM) frequency with $\im\omega<0$\footnote{The definition of quasinormal modes and three equivalent characterizations are given in Chapter~\ref{chap:qnm}.}, the spatial waveforms in Eqs.~\eqref{eq:bh_ingoing_bc} and \eqref{eq:infinity_outgoing_bc} diverge on the real axis. This is not pathological. It is the result of representing, with one complex frequency over all space, a resonance of an open system that decays in time while its wavefront propagates toward infinity. Because of this divergence, however, a QNM waveform is not an ordinary square-integrable $L^2$ eigenfunction. This is the first reason why analytic continuation and the complex scaling introduced in Chapter~\ref{chap:complex-scaling} are needed.

\section{The case of two horizons}

In \term{Schwarzschild--de~Sitter spacetime}, the static region lies between the \term{black-hole horizon} $r_b$ and the \term{cosmological horizon} $r_c$:
\begin{equation}
r_b<r<r_c.
\end{equation}
With a suitable additive constant,
\begin{align}
\rstar&\longrightarrow-\infty,
&&r\longrightarrow r_b+0,
\\
\rstar&\longrightarrow+\infty,
&&r\longrightarrow r_c-0.
\end{align}
The QNM \term{radiation conditions} are an ingoing wave at the black-hole horizon and an outgoing wave at the cosmological horizon. Although the formulas have the same form as Eqs.~\eqref{eq:bh_ingoing_bc} and \eqref{eq:infinity_outgoing_bc}, the physical meaning of the right endpoint is now a cosmological horizon rather than spatial infinity.

By contrast, in an \term{extremal black hole} the zeros of $f(r)$ merge, and the logarithmic behavior in Eq.~\eqref{eq:tortoise_near_horizon} changes. If, near a double root $r=r_e$,
\begin{equation}
f(r)
=
a(r-r_e)^2
+O\bigl((r-r_e)^3\bigr),
\end{equation}
then, with $a$ denoting the quadratic coefficient,
\begin{equation}
\rstar
=
-\frac{1}{a(r-r_e)}
+O\bigl(\log|r-r_e|\bigr).
\end{equation}
Simply substituting extremal parameter values into nonextremal formulas can therefore give the wrong asymptotic structure.

\section{Conclusion of this chapter}

The equation obtained in this chapter is
\begin{equation}
\left(
\partial_t^2
-\partial_{\rstar}^2
+V
\right)\Psi
=
S.
\label{eq:generic_master_equation}
\end{equation}
The type of black hole and the spin of the perturbation enter through the effective potential $V$, while the causal structure enters through the range of $\rstar$ and the asymptotic conditions. This separation turns a geometric problem into a scattering-theory problem. In the next chapter we solve Eq.~\eqref{eq:generic_master_equation} as the response to an external source.

\begin{researchproblem}[Scattering coordinates for an extremal horizon]
For \term{Reissner--Nordstr\"om spacetime}, with $f(r)=1-2M/r+Q^2/r^2$, construct the tortoise coordinates separately in the nonextremal and extremal cases. Exchange the order of the limits $|Q|\to M$ and $r\to r_+$, and determine where the near-horizon wave equation and \term{Jost normalization} become nonuniform. Do not merely reproduce known formulas: propose a uniform coordinate or a \term{matched asymptotic expansion} suitable for later numerical calculations and estimate its error.
\end{researchproblem}

\chapter{Master Fields}
\label{chap:master-field}
\chapterepigraph{\jp{三十輻，共一轂，當其無，有車之用。}}{Laozi, \textit{Tao Te Ching}, Chapter 11}

\section{From geometry to a single field}

The tortoise coordinate maps the exterior region to the real line of a scattering problem. The ten components of a \term{metric perturbation} $h_{\mu\nu}$, however, cannot simply be scattered as they stand. One must separate components removable by coordinate transformations, components fixed by \term{constraint equations}, and components that propagate. A variable---or a small set of variables---that packages the resulting physical content is called a \term{master field}.

We first see the mechanism for a massless scalar field. Take the action to be
\begin{equation}
I[\Phi]
=
-\frac{1}{2}
\int \rmd^4x\sqrt{-g}\,
g^{\mu\nu}\partial_\mu\Phi\partial_\nu\Phi
+\int \rmd^4x\sqrt{-g}\,J\Phi,
\label{eq:scalar-action}
\end{equation}
where $J$ is an external source. Expanding in spherical harmonics and setting $\Phi=r^{-1}\Psi_{\ell m}Y_{\ell m}$, the action after angular integration is, up to boundary terms,
\begin{equation}
I_{\ell m}
=
\frac{1}{2}
\int \rmd t\rmd \rstar
\left[
(\partial_t\Psi_{\ell m})^2
-(\partial_{\rstar}\Psi_{\ell m})^2
-V_\ell\Psi_{\ell m}^2
+2S_{\ell m}\Psi_{\ell m}
\right].
\label{eq:reduced-master-action}
\end{equation}
Variation gives
\begin{equation}
\left(
\partial_t^2-\partial_{\rstar}^2+V_\ell
\right)\Psi_{\ell m}
=
S_{\ell m}.
\label{eq:master-action-eom}
\end{equation}
We retain not only the one-dimensional equation but also the action because the same field will later be quantized.

\section{The minimum needed for gravitational perturbations}

Let $g_{\mu\nu}$ be the background metric and perturb it as
\begin{equation}
g_{\mu\nu}
\longrightarrow
g_{\mu\nu}+h_{\mu\nu}.
\end{equation}
Under an \term{infinitesimal coordinate transformation} $x^\mu\to x^\mu+\xi^\mu$,
\begin{equation}
h_{\mu\nu}
\longrightarrow
h_{\mu\nu}
-\nabla_\mu\xi_\nu
-\nabla_\nu\xi_\mu.
\label{eq:linear-gauge-transformation}
\end{equation}
Individual components of $h_{\mu\nu}$ therefore cannot by themselves be interpreted as physical response. On a Schwarzschild background, decomposition by \term{parity} on the sphere separates the perturbation into \term{axial} and \term{polar} sectors. After the constraints are solved, one \term{gauge-invariant} master field remains in each sector for every $\ell\geq2$ \cite{Regge:1957td,Zerilli:1970wzz,Moncrief:1974am}.

Under inversion on the sphere, $(\vartheta,\varphi)\mapsto(\pi-\vartheta,\varphi+\pi)$, the \term{polar} (even-parity) components transform as $(-1)^\ell$, while the \term{axial} (odd-parity) components transform as $(-1)^{\ell+1}$. One \term{radiative degree of freedom} appears in each sector only for $\ell\geq2$. The $\ell=0$ polar perturbation represents an infinitesimal change of mass plus gauge, while the $\ell=1$ axial perturbation represents an infinitesimal change of angular momentum---that is, the linearized Kerr direction---plus gauge. The $\ell=1$ polar vacuum perturbation is a gauge mode corresponding to a displacement of the center of mass. In the presence of matter sources, constrained \term{monopole and dipole} fields remain, but they are not \term{gravitational radiation} propagating to infinity \cite{Martel:2005ir}. Thus the restriction $\ell\geq2$ in the \term{Regge--Wheeler--Zerilli wave equations} below reflects the existence of radiative degrees of freedom, not a convenience of the potentials.

For the axial \term{Regge--Wheeler field} $\Psi_{\mathrm{RW}}$,
\begin{equation}
\left(
\partial_t^2-\partial_{\rstar}^2+V_{\mathrm{RW}}
\right)\Psi_{\mathrm{RW}}
=
S_{\mathrm{RW}},
\label{eq:rw-master}
\end{equation}
where
\begin{equation}
V_{\mathrm{RW}}(r)
=
f(r)
\left[
\frac{\ell(\ell+1)}{r^2}
-\frac{6M}{r^3}
\right],
\qquad
f(r)=1-\frac{2M}{r}.
\label{eq:rw-potential}
\end{equation}
The polar \term{Zerilli field} satisfies an equation of the same form. Defining
\begin{equation}
\lambda
=
\frac{(\ell-1)(\ell+2)}{2},
\end{equation}
its potential is
\begin{align}
V_{\mathrm Z}(r)
&=
\frac{2f(r)}{r^3(\lambda r+3M)^2}
\Bigl[
\lambda^2(\lambda+1)r^3
+3\lambda^2Mr^2
\nonumber\\
&\hspace{9em}
+9\lambda M^2r
+9M^3
\Bigr].
\label{eq:zerilli-potential-minimum}
\end{align}
The two potentials are different, but they are related by an appropriate first-order differential transformation and have the same quasinormal-mode frequencies, to be defined in Chapter~\ref{chap:qnm} \cite{Chandrasekhar:book,Chandrasekhar:1975nkd}.

\begin{principlebox}
A master field and its effective potential are representations and are not unique. What should be compared across representations are the radiation conditions, the real-frequency \term{scattering matrix}, and the map from physical sources to observables.
\end{principlebox}

\section{Sources and observables}

Vacuum equations alone do not constitute a response theory. If the matter \term{stress tensor} is $T_{\mu\nu}$, spherical-harmonic components of the \term{linearized Einstein equations} determine $S_{\mathrm{RW}}$ or $S_{\mathrm Z}$. Even for the same physical matter source, changing the definition of the master field changes the form of $S$. The master-field solution itself must therefore not be identified directly with an observable.

Let $O$ denote an \term{observation operation}, such as a far-zone radiative waveform, a flux crossing the horizon, or a local curvature perturbation. Schematically, the relation from source to observable is
\begin{equation}
T_{\mu\nu}
\longrightarrow
S_{\ell m}
\longrightarrow
\Psi_{\ell m}
\longrightarrow
\delta O.
\label{eq:source-master-observable-chain}
\end{equation}
The first arrow represents the constraints and projection of the linearized Einstein equations, while the final arrow represents reconstruction of a waveform or flux. In the next chapter, the central arrow will be defined as a solution operator responding to an external source.

\paragraph{One complete chain.}
With the Martel--Poisson normalization, the gravitational wave at future \term{null infinity} can be reconstructed from the polar field $\Psi^{\mathrm{even}}_{\ell m}$ and axial field $\Psi^{\mathrm{odd}}_{\ell m}$ as
\begin{equation}
h_+-\rmi h_\times
=
\frac{1}{2r}
\sum_{\ell m}
\sqrt{\frac{(\ell+2)!}{(\ell-2)!}}
\left(
\Psi^{\mathrm{even}}_{\ell m}
+\rmi\Psi^{\mathrm{odd}}_{\ell m}
\right)
{}_{-2}Y_{\ell m}
+O(r^{-2}).
\label{eq:master-to-strain}
\end{equation}
With $u=t-\rstar$, the radiated energy is
\begin{equation}
\frac{\rmd E}{\rmd u}
=
\frac{1}{64\pi}
\sum_{\ell m}
\frac{(\ell+2)!}{(\ell-2)!}
\left(
\left|\partial_u\Psi^{\mathrm{even}}_{\ell m}\right|^2
+
\left|\partial_u\Psi^{\mathrm{odd}}_{\ell m}\right|^2
\right).
\label{eq:master-to-energy-flux}
\end{equation}
This normalization follows Ref.~\cite{Martel:2005ir}. Other conventions for the RW field insert a time derivative or integral into the odd-parity field and change numerical coefficients. An observable is therefore obtained only by carrying, in one consistent normalization, the projection of the source $T_{\mu\nu}$, the Green-function calculation of $\Psi_{\ell m}$, and the \term{waveform reconstruction} in Eq.~\eqref{eq:master-to-strain}.

\section{What to retain about rotating spacetimes}

Because \term{Kerr spacetime} is not spherically symmetric, it does not reduce to a single Regge--Wheeler-type real potential. Let the mass be $M$ and the angular momentum be $J=aM$, and define
\begin{align}
\Delta
&=
r^2-2Mr+a^2,
&
r_\pm
&=
M\mathord{\pm}\sqrt{M^2-a^2},
&
\Omega_H
&=
\frac{a}{r_+^2+a^2}.
\label{eq:kerr-horizon-minimum}
\end{align}
Here $r_+$ is the event-horizon radius and $\Omega_H$ is its angular velocity. The Kerr tortoise coordinate is defined by
\begin{equation}
\frac{\rmd\rstar}{\rmd r}
=
\frac{r^2+a^2}{\Delta}.
\label{eq:kerr-tortoise-minimum}
\end{equation}

Using curvature components adapted to \term{Petrov type D}, a spin-$s$ \term{Teukolsky variable} can be separated schematically as
\begin{equation}
\psi_s
=
\rme^{-\rmi\omega t+\rmi m\varphi}
S_{s\ell m}(\vartheta;a\omega)
R_{s\ell m}(r).
\label{eq:teukolsky-separation-minimum}
\end{equation}
\cite{Teukolsky:1972my} Here $S_{s\ell m}$ is a \term{spin-weighted spheroidal harmonic}. Because the \term{angular eigenvalue} depends on $a\omega$, the angular and radial equations must be solved together. Suppressing spin-dependent power-law prefactors, the radiation conditions are
\begin{align}
R_{s\ell m}
&\sim
\rme^{-\rmi(\omega-m\Omega_H)\rstar},
&&r\longrightarrow r_+,
\label{eq:kerr-horizon-bc-minimum}
\\
R_{s\ell m}
&\sim
\rme^{+\rmi\omega\rstar},
&&r\longrightarrow\infty.
\label{eq:kerr-infinity-bc-minimum}
\end{align}
The \term{Killing-energy flux} crossing the horizon is proportional to $\omega-m\Omega_H$. Consequently, for bosons in the range $0<\omega<m\Omega_H$, the horizon flux is negative and the reflected wave is stronger than the incident wave: \term{superradiance} occurs. The same frequency shift reappears in the thermal factor in Chapter~\ref{chap:horizon-thermality}.

The minimum amount of Kerr physics retained in the main text is the \term{horizon generator}, \term{Teukolsky separation}, radiation conditions at both ends, and \term{superradiance}. We do not survey metric reconstruction or the full taxonomy of numerical methods; implementation and verification are deferred to the research problem in Chapter~\ref{chap:qnm}.

\section{Conclusion of this chapter}

Black-hole perturbations reduce, after gauge and constraints have been removed, to an \term{open wave problem} for master fields. The action \eqref{eq:reduced-master-action} is the common starting point for classical response and quantization, while Eq.~\eqref{eq:source-master-observable-chain} separates representation-dependent quantities from observables. In the next chapter we invert the master equation in the presence of an external source.

\chapter{The Resolvent}
\label{chap:resolvent}
\chapterepigraph{``Deep calleth unto deep.''}{Psalm 42:7, King James Version}

\section{Ask for the response first}

Abbreviating $x=\rstar$, write the one-dimensional equation obtained in the previous chapter as
\begin{equation}
\left(
\partial_t^2+H
\right)\Psi(t,x)
=
S(t,x).
\label{eq:wave_H_form}
\end{equation}
Here the \term{spatial operator} $H$ is
\begin{equation}
H
=
-\frac{\rmd^2}{\rmd x^2}
+V(x),
\label{eq:H_definition}
\end{equation}
where $V(x)$ is a real-valued effective potential and $S(t,x)$ is an external source.

For a closed system, it is natural to find the \term{eigenfunctions} of $H$ and expand in them. In the black-hole exterior, however, energy escapes through the asymptotic ends on both sides. We therefore change the question.

\begin{principlebox}
Do not begin by asking which complex frequencies are allowed. Ask instead what response is produced by a given source. Then identify poles and the \term{continuous spectrum} as the \term{analytic structure} of that response.
\end{principlebox}

In this order, the quasinormal modes (QNMs) defined in Chapter~\ref{chap:qnm} are not input data but results read off from the response.

\section{Two resolvents}

First introduce a complex energy variable $z\in\bbC$ and define the \term{resolvent} of $H$ by
\begin{equation}
\vR(z)
=
(H-z)^{-1}.
\label{eq:energy_resolvent}
\end{equation}
The set of $z$ for which this inverse exists as a bounded operator is the \term{resolvent set} $\rho(H)$, and its complement is the \term{spectrum} $\sigma(H)$.

Because $z=\omega^2$ for the wave equation, it is also useful to introduce the \term{operator pencil} on the frequency plane,
\begin{equation}
\vL(\omega)
=
H-\omega^2,
\label{eq:operator_pencil}
\end{equation}
and the frequency-domain \term{Green operator}
\begin{equation}
\vG(\omega)
=
\vL(\omega)^{-1}
=
\vR(\omega^2).
\label{eq:frequency_green_operator}
\end{equation}
The $z$ plane makes the spectral theory of a \term{self-adjoint operator} transparent, while the $\omega$ plane makes incoming/outgoing waves and time evolution transparent. We keep the two planes distinct and move between them as needed.

\begin{warningbox}
The map $z=\omega^2$ is not one-to-one. The square root on the $z$ plane has two \term{Riemann sheets}, and the \term{analytic continuation} of radiation conditions includes a choice of \term{sheet}. It is not enough to say that one ``analytically continues across the real axis''; the sign of $\omega$ and the relevant \term{boundary value} must also be specified.
\end{warningbox}

\section{\term{Causal Fourier--Laplace transform}}

Assume that the source acts only after some finite time. Choose $\gamma>0$ sufficiently large and write
\begin{align}
\Psi(t,x)
&=
\int_{\bbR+\rmi\gamma}
\frac{\rmd\omega}{2\pi}
\,
\rme^{-\rmi\omega t}
\psi(\omega,x),
\label{eq:inverse_laplace_psi}
\\
S(t,x)
&=
\int_{\bbR+\rmi\gamma}
\frac{\rmd\omega}{2\pi}
\,
\rme^{-\rmi\omega t}
s(\omega,x).
\label{eq:inverse_laplace_source}
\end{align}
The \term{integration contour} $\bbR+\rmi\gamma$ is a horizontal line in the upper half-plane. Equation~\eqref{eq:wave_H_form} becomes
\begin{equation}
\vL(\omega)\psi(\omega)
=
s(\omega),
\end{equation}
so that
\begin{equation}
\psi(\omega)
=
\vG(\omega)s(\omega).
\label{eq:frequency_response}
\end{equation}
Defining the operator in the upper half-plane and analytically continuing downward selects the \term{boundary value} of the \term{retarded Green function}.

Let the observation operation be a \term{linear functional} $\bra{O}$ and the spatial profile of the source be $\ket{S}$. The observed \term{frequency response} is then
\begin{equation}
\vA(\omega)
=
\bra{O}\vG(\omega)\ket{S}.
\label{eq:source_observer_amplitude}
\end{equation}
The quantity $\vA(\omega)$ is a scalar. Changing the source or the observation point changes the strength with which the same pole appears. The existence of a spectral feature and its strong excitation are separate questions.

For gravitational perturbations, let $\vD_S$ denote the linear map that constructs the master-equation source from the physical matter source $T_{\mu\nu}$, and let $\vD_O$ denote the linear map that constructs the detector observable from the master field. The map from physical source to observed response is then
\begin{equation}
\delta O(\omega)
=
\vD_O\vG(\omega)\vD_S T.
\label{eq:physical-response-map}
\end{equation}
The operators $\vD_S$ and $\vD_O$ depend on the background, the choice of \term{gauge-invariant variables}, and the observation procedure. But if they are transformed consistently, the final $\delta O$ is independent of the representation. This point will be organized systematically in Chapter~\ref{chap:response-invariants}.

\section{Jost solutions}

Suppose that the effective potential approaches zero sufficiently rapidly as $x\to\pm\infty$. Define the left \term{Jost solution} $f_-(x,\omega)$ and right \term{Jost solution} $f_+(x,\omega)$ by
\begin{align}
\vL(\omega)f_-(x,\omega)&=0,
&
f_-(x,\omega)&\sim\rme^{-\rmi\omega x},
&&x\longrightarrow-\infty,
\label{eq:left_jost}
\\
\vL(\omega)f_+(x,\omega)&=0,
&
f_+(x,\omega)&\sim\rme^{+\rmi\omega x},
&&x\longrightarrow+\infty.
\label{eq:right_jost}
\end{align}
If the left endpoint is a black-hole horizon, $f_-$ is the solution entering the horizon; if the right endpoint is spatial infinity, $f_+$ is the solution outgoing to infinity.

Define the \term{Wronskian} of the two solutions by
\begin{equation}
\vW(\omega)
=
f_-(x,\omega)
\frac{\partial f_+(x,\omega)}{\partial x}
-
\frac{\partial f_-(x,\omega)}{\partial x}
f_+(x,\omega).
\label{eq:jost_wronskian}
\end{equation}
Because the differential equation contains no first-derivative term, $\vW(\omega)$ is independent of $x$. Indeed, differentiating Eq.~\eqref{eq:jost_wronskian} with respect to $x$ and using the same differential equation for both Jost solutions gives zero.

\section{Constructing the Green kernel}

Define the frequency-domain \term{Green kernel} $G(x,x';\omega)$ by
\begin{equation}
\vL(\omega)G(x,x';\omega)
=
\delta(x-x').
\label{eq:green_kernel_equation}
\end{equation}
With $x_< =\min(x,x')$ and $x_>=\max(x,x')$, the \term{Green kernel} satisfying the radiation conditions is
\begin{equation}
G(x,x';\omega)
=
-\frac{
f_-(x_<,\omega)f_+(x_>,\omega)
}{
\vW(\omega)
}.
\label{eq:green_jost_formula}
\end{equation}

\paragraph{Matching conditions for the Green kernel.}
For $x\ne x'$, the Green kernel satisfies the homogeneous equation, so it is proportional to $f_-$ on the left and to $f_+$ on the right. Continuity at $x=x'$ reduces the two coefficients to one. Integrating Eq.~\eqref{eq:green_kernel_equation} across a neighborhood of $x'$ further gives
\begin{equation}
\left.
\frac{\partial G}{\partial x}
\right|_{x=x'+0}
-
\left.
\frac{\partial G}{\partial x}
\right|_{x=x'-0}
=
-1.
\end{equation}
This jump in the derivative fixes the sign and normalization in Eq.~\eqref{eq:green_jost_formula}.

Equation~\eqref{eq:green_jost_formula} is one of the central formulas of this book. The numerator represents propagation from the source to the observation point, while the denominator measures whether the radiation conditions at the two ends can be satisfied simultaneously. If $\vW(\omega)=0$, the left and right Jost solutions become linearly dependent, and a nontrivial solution satisfying radiation conditions at both ends exists even without an external source. In the next chapter, such complex frequencies will be defined as quasinormal modes (QNMs).

\section{Scattering coefficients and flux}

Consider a real frequency $\omega>0$. Normalize a scattering solution $\psi_{\mathrm{in}}$ incident from the right by
\begin{align}
\psi_{\mathrm{in}}
&\sim
\rme^{-\rmi\omega x}
+\vS_{\mathrm{R}}(\omega)
\rme^{+\rmi\omega x},
&&x\longrightarrow+\infty,
\label{eq:scatter_right}
\\
\psi_{\mathrm{in}}
&\sim
\vS_{\mathrm{T}}(\omega)
\rme^{-\rmi\omega x},
&&x\longrightarrow-\infty.
\label{eq:scatter_left}
\end{align}
Here $\vS_{\mathrm{R}}$ and $\vS_{\mathrm{T}}$ are the \term{reflection amplitude} and \term{transmission amplitude}, respectively.

Define the \term{flux} associated with a solution $\psi$ by
\begin{equation}
j[\psi]
=
\frac{1}{2\rmi}
\left(
\psi^*\frac{\rmd\psi}{\rmd x}
-
\psi\frac{\rmd\psi^*}{\rmd x}
\right).
\label{eq:conserved_flux}
\end{equation}
For a real potential, $j$ is independent of $x$. If the asymptotic wave numbers on the two sides are equal, \term{flux conservation} gives
\begin{equation}
|\vS_{\mathrm{R}}(\omega)|^2
+
|\vS_{\mathrm{T}}(\omega)|^2
=
1.
\label{eq:unitarity_one_channel}
\end{equation}
The fraction absorbed by the black hole, namely the \term{greybody factor}, is
\begin{equation}
\Gamma(\omega)
=
|\vS_{\mathrm{T}}(\omega)|^2.
\label{eq:greybody_factor}
\end{equation}
When the two asymptotic wave numbers differ, or when several channels are coupled, \term{group velocities} and \term{flux normalization} must also be included.

\section{Analytically continued response}

For a self-adjoint $H$, the physical spectrum lies on the real axis. A \term{resonance} is not added to that spectrum as a new $L^2$ eigenvalue. Instead, it appears as a pole when matrix elements of the \term{Green operator}, initially defined in the upper half-plane, are analytically continued into the lower half-plane while preserving the radiation conditions.

Near a \term{simple pole} $\omega=\omega_n$, one may write
\begin{equation}
\vG(\omega)
=
\frac{A_n}{\omega-\omega_n}
+\vG_{\mathrm{reg}}(\omega).
\label{eq:green_simple_pole}
\end{equation}
Here $A_n$ is the \term{operator-valued residue}, and $\vG_{\mathrm{reg}}$ is the part \term{regular} at $\omega_n$. Sandwiching this expression between the source and observation operation gives
\begin{equation}
\vA(\omega)
=
\frac{
\bra{O}A_n\ket{S}
}{
\omega-\omega_n
}
+\vA_{\mathrm{reg}}(\omega).
\label{eq:amplitude_simple_pole}
\end{equation}
The physically measured strength of the pole is $\bra{O}A_n\ket{S}$, not a QNM waveform with an arbitrarily chosen normalization.

\section{Deforming the contour}

The time-domain response can be analyzed by deforming the \term{integration contour} in Eq.~\eqref{eq:inverse_laplace_psi} downward. Schematically,
\begin{equation}
\Psi(t)
=
\Psi_{\mathrm{pole}}(t)
+\Psi_{\mathrm{cont}}(t)
+\Psi_{\mathrm{arc}}(t).
\label{eq:response_decomposition}
\end{equation}
The first term contains residues of poles crossed during the deformation, the second contains discontinuities across a \term{branch cut} or the continuous spectrum, and the third contains contributions from a \term{large arc} or the high-frequency region.

A simple pole contributes
\begin{equation}
\Psi_{\mathrm{pole}}(t)
\mathrel{\propto}
\rme^{-\rmi\omega_n t}
\bra{O}A_n\ket{S}.
\label{eq:time_simple_pole}
\end{equation}
If $\im\omega_n<0$, this term decays in time. The interval over which Eq.~\eqref{eq:time_simple_pole} dominates, however, is not determined by the pole location alone. One must compare it with the bandwidth of the source, the observation point, the continuous component, and the high-frequency contribution.

\begin{warningbox}
A ``sum over poles'' obtained by contour deformation is not automatically a complete expansion valid at all times. One must check the region in which the large-arc contribution vanishes, the location of the \term{branch cut}, the regularity of the initial data, and the observation point. \term{Completeness of QNMs} must not be inferred by analogy with an eigenvector expansion of a finite-dimensional matrix.
\end{warningbox}

For asymptotically flat Schwarzschild spacetime, \term{nonanalyticity} associated with the long-range potential is related to the late-time power-law component defined in Chapter~\ref{chap:time-response}. For de~Sitter black holes bounded by two horizons, the asymptotic structure is different. Time-domain response is discussed in Chapter~\ref{chap:time-response}, while spectral representations of continuous response are treated in Part II.

\section{Conclusion of this chapter}

The basic object constructed here is not a set of allowed frequencies but the Green operator $\vG(\omega)$. The Jost formula \eqref{eq:green_jost_formula} shows that singularities of scattering amplitudes and of externally driven response are governed by the same \term{Wronskian}. In the next chapter we define its poles as quasinormal modes (QNMs) and study them in detail.

\begin{researchproblem}[Long-range Jost functions]
The Schwarzschild effective potential is not a finite-range potential at infinity. Reconstruct the usual Jost-solution construction based on a \term{Volterra equation} while retaining the logarithmic difference between $\rstar$ and $r$. Identify the location and sheet of the low-frequency \term{branch point}, and design a method that extracts both zeros of the Wronskian and the discontinuity across the branch cut from the same numerical calculation. As a convergence test, require the time-domain response to remain invariant when the placement of the cut is changed.
\end{researchproblem}

\chapter{Quasinormal Modes}
\label{chap:qnm}
\chapterepigraph{``Rest, rest, perturb\`ed spirit.''}{William Shakespeare, \textit{Hamlet}, Act I, Scene 5}

\section{Three equivalent definitions}

The Green kernel obtained in the previous chapter is
\begin{equation}
G(x,x';\omega)
=
-\frac{f_-(x_<,\omega)f_+(x_>,\omega)}{\vW(\omega)}.
\label{eq:qnm-start-green}
\end{equation}
When $\vW(\omega)$ vanishes, the solution satisfying the outgoing condition on the left and the solution satisfying the outgoing condition on the right become the same solution.

\begin{definition}[Quasinormal mode]
A complex frequency $\omega_n$ satisfying
\begin{equation}
\vW(\omega_n)=0
\label{eq:qnm_wronskian_zero}
\end{equation}
is called a \term{QNM frequency}. The corresponding nontrivial homogeneous solution $\psi_n(x)$ satisfies radiation conditions directed out of the system at both asymptotic ends.
\end{definition}

With the time factor $\rme^{-\rmi\omega t}$, the conditions in the black-hole exterior are
\begin{align}
\psi_n(x)
&\sim
\rme^{-\rmi\omega_n x},
&&x\longrightarrow-\infty,
\label{eq:qnm_bc_left}
\\
\psi_n(x)
&\sim
\rme^{+\rmi\omega_n x},
&&x\longrightarrow+\infty.
\label{eq:qnm_bc_right}
\end{align}
Outgoing through the left endpoint means entering the event horizon. Thus a \term{QNM} has three equivalent characterizations: a solution of the \term{radiation boundary-value problem}, a zero of the Jost Wronskian, and a pole of the analytically continued Green function.

\section{Simple poles and residues}

If $\vW'(\omega_n)\ne0$, then $\omega_n$ is a \term{simple pole}, and
\begin{equation}
\mathop{\mathrm{Res}}_{\omega=\omega_n}
G(x,x';\omega)
=
-\frac{
f_-(x_<,\omega_n)f_+(x_>,\omega_n)
}{
\vW'(\omega_n)
}.
\label{eq:qnm_green_residue}
\end{equation}
Multiplying a Jost solution by a constant rescales numerator and denominator by the same factor, so this \term{residue} is independent of normalization.

The \term{QNM} waveform $\psi_n$ itself, however, can be multiplied by an arbitrary constant. Let $\bra{O}$ be the observation operation, $\ket{S}$ the source, and $A_n$ the \term{operator-valued residue}. The observed pole contribution is
\begin{equation}
\vA_n(\omega)
=
\frac{a_n}{\omega-\omega_n},
\qquad
a_n
=
\bra{O}A_n\ket{S}.
\label{eq:qnm-observable-residue}
\end{equation}
The background-determined frequency, a normalization-dependent waveform, the operator-valued \term{residue}, and an amplitude specified by a source and observer are not the same quantity \cite{Motohashi:2024fwt}.

\begin{principlebox}
Frequencies alone do not determine a time-domain waveform. The physical \term{mode strength} is the \term{matrix element} obtained by sandwiching the Green-function residue between the source and the observation operation.
\end{principlebox}

\section{Spatial divergence and temporal decay}

For a damped QNM of a stable black hole, $\im\omega_n<0$. At a fixed position,
\begin{equation}
\left|\rme^{-\rmi\omega_n t}\right|
=
\rme^{(\im\omega_n)t}
\end{equation}
decays. But at the right endpoint the spatial factor satisfies
\begin{equation}
\left|\rme^{+\rmi\omega_n x}\right|
=
\rme^{-(\im\omega_n)x},
\end{equation}
which grows as $x\to+\infty$; the same happens at the left endpoint. This is the result of extending an outward-propagating wavefront over all space with a single complex frequency. It is not a temporal instability.

\begin{warningbox}
A QNM is not an ordinary $L^2$ eigenfunction. If its divergence on the real axis is ignored and the QNMs are treated as a \term{complete set} in a \term{Hilbert space}, both normalization and the domain of a mode expansion can be misidentified. The \term{complex scaling} defined in Chapter~\ref{chap:complex-scaling} converts this divergence into decay along an analytically continued coordinate.
\end{warningbox}

\section{Energy fixes the direction of damping}

For a real potential $V(x)\geq0$ and a real solution of \term{finite energy}, define the energy on $[-R,R]$ by
\begin{equation}
E_R(t)
=
\frac{1}{2}
\int_{-R}^{R}\rmd x
\left[
(\partial_t\Psi)^2
+(\partial_x\Psi)^2
+V\Psi^2
\right].
\label{eq:qnm-energy-estimate}
\end{equation}
Using the wave equation and integrating by parts gives
\begin{equation}
\frac{\rmd E_R}{\rmd t}
=
\left[\partial_t\Psi\,\partial_x\Psi\right]_{-R}^{R}.
\label{eq:qnm-energy-flux-balance}
\end{equation}
If only right-moving waves are allowed at the right endpoint and left-moving waves at the left endpoint, then as $R\to\infty$ the right-hand side is the negative flux escaping through the two ends, and the exterior energy cannot increase.

Suppose there were an outgoing mode with $\im\omega>0$. Its spatial profile would decay at both ends and have finite energy, while its time factor would grow exponentially. This contradicts Eq.~\eqref{eq:qnm-energy-flux-balance}. Hence this positive-energy problem has no \term{unstable QNMs}. The argument does not determine the damping rate and cannot be applied unchanged when $V$ is not positive, when rotation permits negative horizon flux, or when gravitational constraints must be included. Stability is not an empirical pattern in a frequency table; it is a question of proving the sign of an appropriate energy and its \term{boundary flux}.

\section{A solvable barrier}

To see the origin of QNMs in one formula, consider the \term{P\"oschl--Teller barrier}
\begin{equation}
V(x)
=
\frac{V_0}{\cosh^2(\alpha x)},
\qquad
V_0>0,
\quad
\alpha>0.
\label{eq:poschl_teller_potential}
\end{equation}
With the change of variable $y=(1+\tanh\alpha x)/2$, the equation becomes a \term{hypergeometric equation}. Imposing the radiation conditions at both ends gives
\begin{equation}
\omega_n
=
\mathord{\pm}
\sqrt{V_0-\frac{\alpha^2}{4}}
-\rmi\alpha\left(n+\frac{1}{2}\right),
\qquad
n=0,1,2,\ldots.
\label{eq:poschl_teller_qnm}
\end{equation}
The barrier height controls the real part, while its width controls the spacing of damping rates. This is not the exact black-hole potential, but it shows that open boundary conditions alone generate complex frequencies.

\section{The photon-sphere limit}

For Schwarzschild spacetime with $\ell\gg1$, the leading part of the effective potential is $f(r)\ell(\ell+1)/r^2$, whose maximum lies at the \term{unstable photon orbit} $r=3M$. The angular frequency and \term{Lyapunov exponent} of that orbit are both
\begin{equation}
\Omega_c
=
\Lambda_c
=
\frac{1}{3\sqrt{3}M}.
\end{equation}
For fixed \term{overtone} number $n$,
\begin{equation}
\omega_{\ell n}
\simeq
\Omega_c\left(\ell+\frac{1}{2}\right)
-\rmi\Lambda_c\left(n+\frac{1}{2}\right).
\label{eq:eikonal_qnm}
\end{equation}
Local ray dynamics explains the slowly damped part of the spectrum, but finite-$\ell$ poles, residues, and the \term{branch cut} are not obtained from this approximation. For example, the fundamental $\ell=2$ gravitational mode is
\begin{equation}
M\omega_{20}
\simeq
0.37367-0.08896\rmi.
\label{eq:schwarzschild_fundamental_qnm}
\end{equation}
\cite{Berti:2009kk}

\section{Diagnostics for numerical calculations}

Direct integration finds zeros of the Wronskian of the left and right Jost solutions. The \term{continued-fraction method} imposes a \term{minimal-solution condition} on an \term{asymptotic series}. Time evolution estimates poles from the waveform. The complex-scaling method introduced in Chapter~\ref{chap:complex-scaling} converts outgoing waves into decaying boundary conditions. Although the methods differ, their diagnostics are common.

\begin{enumerate}[label=(\roman*),leftmargin=2.8em]
\item The time convention and the radiation conditions at both ends must be consistent.
\item The poles must remain stable when the basis size, domain, and numerical precision are varied.
\item Discretized continuous components must be separated from genuine poles.
\item If a physical amplitude is required, the residue must converge as well.
\end{enumerate}

Complex-scaling calculations for Schwarzschild and Reissner--Nordstr\"om are given in Ref.~\cite{Ogawa:2026veu}. In a \term{non-self-adjoint problem}, sensitivity of eigenvalues must be distinguished from sensitivity of the real-frequency response \cite{Jaramillo:2020tuu}.

\section{Conclusion of this chapter}

A QNM is not a material vibration intrinsic to a black hole. It is a resonance pole of a Green function with open boundary conditions. The pole is important, but a response also requires the residue, source, observation operation, and continuous components. In the next chapter we invert the pole and continuous contributions and ask which parts of the time-domain waveform they control.

\begin{researchproblem}[Numerical implementation of Kerr QNMs]
Solve the angular and radial Teukolsky equations simultaneously and track the fundamental QNM with $s=-2$ and $\ell=m=2$ from $a/M=0$ into the near-extremal regime \cite{Teukolsky:1972my,Leaver:1985ax}. At minimum, construct an implementation satisfying the following requirements.

\begin{enumerate}[label=(\roman*),leftmargin=2.8em]
\item Determine the angular eigenvalue and $\omega$ in the same iteration or \term{contour method}.
\item Vary the series order, numerical domain, and precision to separate stable poles from discretized continuous components.
\item Verify that the Schwarzschild value is recovered as $a\to0$.
\item If possible, compute the Green-function residue and identify quantities invariant under a change of normalization.
\end{enumerate}

Reproducing a frequency table is only the entry point. Formulate, from the implementation itself, a research problem concerning at least one of the following: the near-extremal \term{branch structure}, the \term{zero-damped mode}, or the source-dependent \term{excitation amplitude}.
\end{researchproblem}

\chapter{Time-Domain Response}
\label{chap:time-response}
\chapterepigraph{``To every thing there is a season, and a time to every purpose under the heaven.''}{Ecclesiastes 3:1, King James Version}

\section{What is ringdown?}

When a binary black hole merges, the observed gravitational wave passes in sequence through the growing-amplitude binary motion, a strongly nonlinear merger, and a decaying post-merger waveform (Fig.~\ref{fig:ringdown-regimes}). We call \term{ringdown} the time interval in which the post-merger remnant can be approximated as a single stationary black hole plus a small perturbation. This is not synonymous with ``the part of the waveform after its maximum amplitude.'' Immediately after the peak, nonlinear merger effects can remain, while at sufficiently late times a \term{power-law decay} generated by continuous components can dominate.

In a regime dominated by isolated quasinormal-mode poles, the dimensionless \term{gravitational-wave strain} $h(t)$ can be represented schematically as
\begin{align}
h_{\mathrm{RD}}(t)
&=
2\re\sum_{n\in C}
A_n\rme^{-\rmi\omega_n(t-t_0)},
\label{eq:ringdown-waveform-model}
\\
\omega_n
&=
\Omega_n-\frac{\rmi}{\tau_n},
\qquad
\Omega_n>0,
\quad
\tau_n>0.
\label{eq:ringdown-frequency-decay}
\end{align}
Here $A_n$ is a complex amplitude containing both the source and the observation operation, $\Omega_n$ is the angular oscillation frequency, $\tau_n$ is the damping time, and $t_0$ is a reference time chosen to parametrize phase and amplitude. Equation~\eqref{eq:ringdown-waveform-model} is not the definition of ringdown; it is a \term{pole approximation} valid in part of the ringdown regime.

\begin{figure}[tbp]
\centering
\resizebox{\linewidth}{!}{\begin{tikzpicture}[
  scale=1.1,
  every node/.style={font=\scriptsize},
  >=Latex
]
% Binary, merger, and post-merger black hole
\fill[black] (1.00,2.45) circle (.15);
\fill[black] (1.55,2.45) circle (.15);
\draw[gray,-{Latex[length=2mm]},line width=.45pt]
  (1.28,2.74)
  arc[start angle=90,end angle=360,x radius=.47,y radius=.28];
\node[gray] at (1.28,2.04) {binary};

\draw[gray,-{Latex[length=2mm]}]
  (2.00,2.45)--(2.72,2.45);

\fill[black] (3.16,2.45) circle (.23);
\fill[black] (3.46,2.45) circle (.23);
\draw[red,line width=.55pt]
  (3.00,2.45)
  .. controls (3.18,2.83) and (3.50,2.83) ..
  (3.66,2.45);
\node[red] at (3.31,2.04) {merger};

\draw[gray,-{Latex[length=2mm]}]
  (3.86,2.45)--(4.58,2.45);

\fill[black] (5.02,2.45) circle (.25);
\draw[blue,line width=.65pt]
  (4.64,2.45)
  arc[start angle=180,end angle=0,x radius=.38,y radius=.19];
\draw[blue!65,line width=.45pt]
  (4.56,2.45)
  arc[start angle=180,end angle=0,x radius=.46,y radius=.28];
\node[blue] at (5.02,2.04) {post-merger};

% Domains of validity
\fill[red!7]
  (3.45,-.92) rectangle (4.45,1.12);
\fill[blue!5]
  (4.45,-.92) rectangle (11.15,1.12);

\draw[red,line width=1.4pt]
  (3.45,1.31)--(4.45,1.31);
\node[red,above] at (3.95,1.31)
  {nonlinear general relativity};

\draw[blue,line width=1.4pt]
  (4.45,1.31)--(11.15,1.31);
\node[blue,above] at (7.80,1.31)
  {linear black-hole perturbation theory};
\node[blue,below,align=center] at (7.80,1.31)
  {$g_{\mu\nu}=g^{\mathrm{BH}}_{\mu\nu}+\delta g_{\mu\nu}$,
   $|\delta g|\ll|g^{\mathrm{BH}}|$};

% Waveform axes
\draw[gray,-{Latex[length=2mm]}]
  (.18,0)--(11.42,0)
  node[right] {$t$};
\draw[gray,-{Latex[length=2mm]}]
  (.34,-.92)--(.34,1.02)
  node[above] {$h(t)$};

\def\tpeak{3.48}
\def\tlin{4.45}
\def\ttail{8.10}

% Match amplitude and phase at the joining points.
\pgfmathsetmacro{\Apeak}{.055*exp(.60*\tpeak)}
\pgfmathsetmacro{\phipeak}{74*\tpeak+25*\tpeak*\tpeak}
\pgfmathsetmacro{\Alin}{\Apeak*exp(-.10*(\tlin-\tpeak))}
\pgfmathsetmacro{\philin}{\phipeak+395*(\tlin-\tpeak)}
\pgfmathsetmacro{\Atail}{
  \Alin
  *exp(-.50*(\ttail-\tlin))
  *sin(\philin+420*(\ttail-\tlin))
}

% Inspiral
\draw[
  black,line width=.75pt,smooth,
  domain=.45:\tpeak,samples=220
]
  plot
  (\x,{
    .055*exp(.60*\x)
    *sin(74*\x+25*\x*\x)
  });

% Nonlinear merger
\draw[
  black,line width=.95pt,smooth,
  domain=\tpeak:\tlin,samples=100
]
  plot
  (\x,{
    \Apeak
    *exp(-.10*(\x-\tpeak))
    *sin(\phipeak+395*(\x-\tpeak))
  });

% Exponential envelope of the QNM ringdown
\draw[
  red!60,dashed,line width=.45pt,smooth,
  domain=\tlin:\ttail,samples=80
]
  plot
  (\x,{
    \Alin*exp(-.50*(\x-\tlin))
  });
\draw[
  red!60,dashed,line width=.45pt,smooth,
  domain=\tlin:\ttail,samples=80
]
  plot
  (\x,{
    -\Alin*exp(-.50*(\x-\tlin))
  });

% QNM-dominated ringdown
\draw[
  red,line width=1.05pt,smooth,
  domain=\tlin:\ttail,samples=220
]
  plot
  (\x,{
    \Alin
    *exp(-.50*(\x-\tlin))
    *sin(\philin+420*(\x-\tlin))
  });

% Late-time power-law tail
\draw[
  blue,line width=1.05pt,smooth,
  domain=\ttail:11.15,samples=100
]
  plot
  (\x,{
    \Atail/pow(1+.62*(\x-\ttail),2.4)
  });

% Transition times
\draw[gray,densely dotted]
  (3.48,-1.02)--(3.48,1.05);
\draw[blue,densely dashed]
  (4.45,-1.02)--(4.45,1.12);
\draw[gray,densely dotted]
  (8.10,-1.02)--(8.10,1.05);

\node[below,gray] at (3.48,-1.00)
  {$t_{\mathrm{peak}}$};
\node[below,blue] at (4.45,-1.00)
  {$t_{\mathrm{lin}}$};
\node[below,gray] at (8.10,-1.00)
  {$t_{\mathrm{tail}}$};

% Ringdown and tail
\draw[red,line width=1.0pt]
  (4.58,-1.46)--(7.96,-1.46);
\node[below,red] at (6.24,-1.46)
  {QNM-dominated ringdown};

\draw[blue,line width=1.0pt]
  (8.24,-1.46)--(11.03,-1.46);
\node[below,blue,align=center] at (9.64,-1.46)
  {tail: $h\sim t^{-p}$, $p>0$\\
   branch cut / continuous component};
\end{tikzpicture}
}
\caption{Schematic waveform from binary coalescence to late times. $t_{\mathrm{peak}}$ lies near the maximum amplitude, $t_{\mathrm{lin}}$ is the time at which linear perturbation theory becomes accurate to a specified tolerance, and $t_{\mathrm{tail}}$ is the time at which the power-law decay from the continuous component can no longer be neglected. The latter two are not universal boundaries; they depend on the source, observable, and required accuracy.}
\label{fig:ringdown-regimes}
\end{figure}

It is therefore not generally correct either to fix the start of ringdown at $t_{\mathrm{peak}}$ or to assume that the waveform is exactly represented by a finite number of QNMs. Below we define this time-domain regime through the inverse transform of the causal Green function and ask when the \term{pole approximation} is justified.

\section{Moving the initial-value problem to frequency space}

Let the master-field equation be
\begin{equation}
\left(\partial_t^2+H\right)\Psi(t)
=
S(t).
\label{eq:time-response-wave}
\end{equation}
Denote the \term{initial data} at $t=0$ by $\Psi_0$ and $\Pi_0=\partial_t\Psi(0)$. Using the \term{one-sided Fourier--Laplace transform}
\begin{equation}
\widehat\Psi(\omega)
=
\int_0^\infty \rmd t\,
\rme^{+\rmi\omega t}\Psi(t),
\qquad
\im\omega>0,
\label{eq:one-sided-transform}
\end{equation}
integration by parts gives
\begin{equation}
\left(H-\omega^2\right)\widehat\Psi(\omega)
=
\widehat S(\omega)
+\Pi_0
-\rmi\omega\Psi_0.
\label{eq:initial-data-source}
\end{equation}
Thus the \term{initial data} enter the Green operator as an \term{effective source}:
\begin{equation}
\widehat\Psi(\omega)
=
\vG(\omega)
\left[
\widehat S(\omega)+\Pi_0-\rmi\omega\Psi_0
\right].
\label{eq:initial-response-resolvent}
\end{equation}
This single line already shows that QNM excitation depends on the \term{initial data}.

\section{The contour separates the waveform}

For $t>0$, the field is recovered from an upper-half-plane contour $C_+$ as
\begin{equation}
\Psi(t)
=
\int_{C_+}
\frac{\rmd\omega}{2\pi}\,
\rme^{-\rmi\omega t}
\widehat\Psi(\omega).
\label{eq:time-inverse-contour}
\end{equation}
Deforming the contour downward gives, schematically,
\begin{equation}
\Psi(t)
=
\Psi_{\mathrm{arc}}(t)
+\Psi_{\mathrm{QNM}}(t)
+\Psi_{\mathrm{cut}}(t).
\label{eq:prompt-ringdown-tail}
\end{equation}
The large-frequency arc gives the \term{prompt response} immediately behind the wavefront; isolated poles give QNM-dominated ringdown; and the \term{discontinuity} across a \term{branch cut} gives the \term{late-time tail}. These are not separate solutions but different pieces of one contour representation \cite{Leaver:1986gd}.

The sum over simple poles is
\begin{equation}
\Psi_{\mathrm{QNM}}(t)
=
-\rmi
\sum_{n\in C}
\rme^{-\rmi\omega_n t}
A_n
\left[
\widehat S(\omega_n)+\Pi_0-\rmi\omega_n\Psi_0
\right].
\label{eq:qnm-time-sum}
\end{equation}
Here $C$ denotes the set of poles crossed during the deformation. The overall sign depends on the contour orientation, but the pairing of pole exponentials and residues does not.

\section{The domain of a QNM expansion}

Equation~\eqref{eq:qnm-time-sum} is not a \term{complete expansion} valid at all times. Keeping only the QNM sum requires the large-arc contribution to vanish for the observation time and position, the \term{branch-cut} contribution to be negligible, and the pole sum to converge. In particular, summing QNMs before the wavefront arrives cannot by itself reproduce \term{causality}.

At very late times, on the other hand, a branch point near the origin can generate a \term{power-law decay} that dominates even the pole closest to the real axis in the lower half-plane. Ringdown is therefore generally an \term{intermediate-time asymptotic regime}.

\begin{warningbox}
The fact that a waveform can be fitted by finitely many damped sinusoids is not the same as saying that the Green function is represented by those poles alone. Vary the start and end times of the fit, the treatment of continuous components, and the \term{noise model}, and test the stability of pole locations and amplitudes separately.
\end{warningbox}

\section{Discontinuity across a branch cut}

Let $\vG_+$ and $\vG_-$ be the boundary values on the two sides of a branch cut, and define their \term{discontinuity} by
\begin{equation}
\operatorname{Disc}\vG(\omega)
=
\vG_+(\omega)-\vG_-(\omega).
\label{eq:green-discontinuity}
\end{equation}
Writing the branch cut as $C_{\mathrm{cut}}$, the continuous contribution is
\begin{equation}
\Psi_{\mathrm{cut}}(t)
=
\int_{C_{\mathrm{cut}}}
\frac{\rmd\omega}{2\pi}\,
\rme^{-\rmi\omega t}
\operatorname{Disc}\vG(\omega)
\ket{S_{\mathrm{eff}}(\omega)}.
\label{eq:cut-time-response}
\end{equation}
Here $\ket{S_{\mathrm{eff}}}$ collects the terms on the right-hand side of Eq.~\eqref{eq:initial-data-source}. The late-time power law is determined by the \term{small-frequency expansion} of the discontinuity near the origin.

In a spacetime bounded by two horizons, the potential decays exponentially at both ends. One must therefore not assume the same branch-cut structure as in an asymptotically flat spacetime. The boundary conditions must be used anew to determine which \term{nonanalyticities} are physical and which point sequences are artifacts of \term{numerical discretization}.

\section{From source to observer}

Let $A(t)$ be the detector signal and write its frequency-domain response as
\begin{equation}
\vA(\omega)
=
\bra{O}\vG(\omega)\ket{S}.
\label{eq:observable-frequency-response}
\end{equation}
Pole positions do not depend on the source or observer, but whether a given pole is visible depends on $\bra{O}A_n\ket{S}$. This matrix element may vanish by symmetry, nearby poles may cancel each other, or a continuous contribution may mask a pole. The ``\term{mode content}'' of ringdown is not a property of the background alone.

\begin{principlebox}
A decomposition of the time-domain waveform is a contour decomposition of the frequency-domain Green function. The boundaries between \term{prompt response}, ringdown, and the \term{late-time tail} are not absolute, but the analytic structures that generate them can be distinguished.
\end{principlebox}

\section{Conclusion of this chapter}

Initial data and external forcing enter the same Green operator, and poles, branch cuts, and the high-frequency arc combine into one causal time-domain response. This closes the classical-field response theory developed in Part I. In Part II we move poles and continuous components that are awkward to handle on the real axis onto spectra defined along complex coordinates.

\begin{researchproblem}[A uniform approximation for poles and tails]
For a fixed $\ell$ in Schwarzschild spacetime, numerically evaluate both the sum over QNM poles and the low-frequency branch-cut integral from the same \term{initial data}. Construct a \term{uniform approximation} that predicts the crossover time between ringdown and the tail using not only the smallest damping rate but also the residues and the support of the initial data. The diagnostic criterion is that a single error norm should control the discrepancy from direct time evolution after the wavefront arrives.
\end{researchproblem}


\part{The Spectrum of Resonances}
\chapter{Complex Scaling}
\label{chap:complex-scaling}
\chapterepigraph{\jp{大音希聲，大象無形。}}{Laozi, \textit{Tao Te Ching}, Chapter 41}

\section{Turning outgoing waves into decaying waves}

In this chapter we define the \term{complex scaling method} (\term{CSM}), which deforms the coordinate into the complex plane so that resonant waves become \term{square integrable}. Consider a QNM waveform at the right endpoint,
\begin{equation}
\psi(x)
\sim
\rme^{+\rmi\omega x},
\qquad
\omega=\omega_R-\rmi\gamma,
\qquad
\gamma>0.
\end{equation}
On the real axis, $|\psi|\sim\rme^{\gamma x}$ grows. We therefore rotate the coordinate according to
\begin{equation}
x
=
y\rme^{\rmi\theta},
\qquad
0<\theta<\frac{\pi}{2},
\label{eq:uniform-complex-scaling}
\end{equation}
where $y>0$ is real. The magnitude of the wave becomes
\begin{equation}
\left|
\rme^{+\rmi\omega y\rme^{\rmi\theta}}
\right|
=
\exp\left[
-y\left(
\omega_R\sin\theta-\gamma\cos\theta
\right)
\right].
\label{eq:complex-scaling-decay}
\end{equation}
Thus, if $\omega_R\sin\theta>\gamma\cos\theta$, a resonant wave that diverges on the real axis decays along the \term{complex contour}.

For a black hole, both asymptotic ends are rotated outward. One example of \term{exterior complex scaling}, which leaves the interaction region $|y|\leq R_0$ on the real axis, is
\begin{equation}
x_\theta(y)
=
\begin{cases}
y,
& |y|\leq R_0,
\\
\sgn(y)
\left[
R_0+(|y|-R_0)\rme^{\rmi\theta}
\right],
& |y|>R_0.
\end{cases}
\label{eq:exterior-complex-scaling}
\end{equation}

\section{The rotated operator}

For a \term{uniform rotation}, $\rmd/\rmd x=\rme^{-\rmi\theta}\rmd/\rmd y$, and hence
\begin{equation}
H_\theta
=
-\rme^{-2\rmi\theta}
\frac{\rmd^2}{\rmd y^2}
+V\left(y\rme^{\rmi\theta}\right).
\label{eq:complex-scaled-hamiltonian}
\end{equation}
The operator $H_\theta$ is \term{non-self-adjoint}. In the energy plane $z$, the \term{continuous spectrum} of the free operator rotates as
\begin{equation}
[0,\infty)
\longrightarrow
\rme^{-2\rmi\theta}[0,\infty).
\label{eq:rotated-continuum}
\end{equation}
Resonances exposed between the rotated ray and the real axis become \term{discrete eigenvalues} independent of $\theta$, within the domain permitted by analyticity.

\begin{warningbox}
Complex scaling requires a region in which the potential can be continued onto the \term{complex contour}. The admissible contours near a horizon, at infinity, and near an extremal horizon need not be the same. Stationarity of an eigenvalue as $\theta$ is varied is a diagnostic, not a substitute for proving analyticity.
\end{warningbox}

A construction for Schwarzschild and Reissner--Nordstr\"om QNMs is given in Ref.~\cite{Ogawa:2026veu}. Because the radiation conditions are converted into a matrix eigenvalue problem, pole searches and regularization of the resonant waveforms can be carried out within the same computation.

\section{Discrete poles and \term{discretized continuous components}}

Diagonalizing $H_\theta$ in a \term{finite basis} produces isolated resonances together with a sequence of eigenvalues along the rotated ray. Let the \term{right and left eigenvectors} be $\ket{\phi^R_\nu}$ and $\bra{\phi^L_\nu}$, and impose the \term{biorthogonal normalization}
\begin{equation}
\langle\phi^L_\mu\mid\phi^R_\nu\rangle
=
\delta_{\mu\nu}.
\end{equation}
Then the finite-basis resolvent is approximately
\begin{equation}
\vR_\theta(z)
\simeq
\sum_\nu
\frac{
\ket{\phi^R_\nu}\bra{\phi^L_\nu}
}{
z_\nu-z
}.
\label{eq:csm-finite-resolvent}
\end{equation}
The sequence that rotates with $\theta$ is not a numerical error. Its points act as \term{numerical quadrature points} for the \term{continuous-spectrum} integral and carry the non-pole part of the Green function.

\begin{principlebox}
Not every \term{complex eigenvalue} produced by complex scaling is a QNM. Isolated points that remain stationary as the \term{rotation angle} and basis size are varied are resonances; sequences lying along the rotated ray are a discrete representation of the continuous response. Both are needed to reconstruct the response.
\end{principlebox}

\section{Continuum level density}

Let $H_\theta$ be the interacting operator and $H_{0,\theta}$ a \term{comparison operator}. For real energy $E$, define from their \term{resolvent difference}
\begin{equation}
\Delta\rho(E)
=
\frac{1}{\pi}
\im\Tr\left[
\vR_\theta(E+\rmi0)
-\vR_{0,\theta}(E+\rmi0)
\right].
\label{eq:continuum-level-density}
\end{equation}
Even if the individual \term{traces} diverge, their difference can be finite under appropriate conditions. The quantity $\Delta\rho$ maps back onto the real axis the continuous response that is invisible if one looks only at isolated poles \cite{Suzuki:2005wv}.

In Schwarzschild--de~Sitter spacetime, both ends of the exterior region are horizons and the potential decays exponentially at each end. QNMs and continuous components of scalar, electromagnetic, and gravitational perturbations can be treated within the same complex-scaled representation \cite{Ogawa:2026dSCSM}. Tracking how not only the poles but also $\Delta\rho$ move as the cosmological constant is varied gives a description of changes in the time-domain response that contains more information than a list of poles.

\section{Three checks to retain in numerical work}

A CSM calculation should preserve at least the following diagnostics.

\begin{enumerate}[label=(\roman*),leftmargin=2.8em]
\item Eigenvalue plots for several values of the \term{rotation angle} $\theta$.
\item Convergence of poles, residues, and the \term{continuum level density} as the basis size is varied.
\item Comparison with a Green function computed directly at real frequencies.
\end{enumerate}

A physical conclusion should not be drawn from a single complex-eigenvalue plot. The third check verifies that the representation on the complex contour reproduces the original scattering response.

\section{Conclusion of this chapter}

Complex scaling converts divergent QNMs into decaying eigenfunctions and places isolated poles and continuous components within a single \term{non-self-adjoint spectrum}. In the next chapter we study what happens when discrete poles themselves coalesce and an individual-mode representation breaks down.

\begin{researchproblem}[Universality of continuous response in de~Sitter systems]
For several spins or \term{coupled channels} in Schwarzschild--de~Sitter and Reissner--Nordstr\"om--de~Sitter spacetimes, implement the same exterior complex scaling and \term{trace subtraction}. Vary the cosmological constant, charge, and dimension, and determine which combination of poles and $\Delta\rho(E)$ yields a uniform approximation to the time-domain response. The final diagnostic should be convergence of the sourced real-axis response, not the arrangement of the $\theta$-dependent point sequence.
\end{researchproblem}

\chapter{Exceptional Points}
\label{chap:exceptional-points}
\chapterepigraph{``The way up and the way down are one and the same.''}{Heraclitus, fragment B60 (Diels--Kranz)}

\section{When two resonances become one}

Suppose that the background or effective potential depends on two real control parameters $\lambda_1$ and $\lambda_2$. A point at which two resonance eigenvalues coincide and their right eigenvectors also coalesce into one is called a second-order \term{exceptional point}. Unlike an ordinary \term{degeneracy}, the dimension of the \term{eigenspace} decreases.

In this chapter we call the subspace containing these two resonances, isolated from the rest of the spectrum, a \term{cluster}. The operator restricted to the cluster subspace can be written as
\begin{equation}
H_{\mathrm{cl}}
=
z_{\mathrm{EP}}P+N,
\qquad
N^2=0,
\qquad
N\ne0,
\label{eq:ep-jordan-form}
\end{equation}
where $P$ is the \term{projection} onto the two-dimensional cluster and $N$ is the \term{nilpotent part}. If a single complexified control parameter $\lambda$ is used, the nearby eigenvalues generically take the form
\begin{equation}
z_\pm(\lambda)
=
z_{\mathrm{EP}}
\mathord{\pm}
c\sqrt{\lambda-\lambda_{\mathrm{EP}}}
+O(\lambda-\lambda_{\mathrm{EP}}),
\label{eq:ep-puiseux}
\end{equation}
with a system-dependent constant $c\ne0$. Encircling the exceptional point exchanges the two \term{sheets} of the \term{square root} and swaps the two states.

\section{The minimal two-by-two model}

To see the \term{square-root branch point} and the \term{double pole} simultaneously, consider
\begin{equation}
H(\lambda)
=
\begin{pmatrix}
0&1\\
\lambda&0
\end{pmatrix}.
\label{eq:ep-two-by-two-model}
\end{equation}
For $\lambda\ne0$, the eigenvalues and right eigenvectors are
\begin{equation}
\omega_\pm
=
\mathord{\pm}\sqrt{\lambda},
\qquad
v_\pm
=
\begin{pmatrix}
1\\
\omega_\pm
\end{pmatrix}.
\end{equation}
As $\lambda\to0$, not only the two eigenvalues but also the two eigenvectors coalesce into $(1,0)^{\mathrm T}$.

Direct inversion of the \term{frequency resolvent} gives
\begin{equation}
\left(\omega I-H(\lambda)\right)^{-1}
=
\frac{1}{\omega^2-\lambda}
\begin{pmatrix}
\omega&1\\
\lambda&\omega
\end{pmatrix},
\label{eq:ep-two-by-two-resolvent}
\end{equation}
and at the exceptional point $\lambda=0$,
\begin{equation}
\left(\omega I-H(0)\right)^{-1}
=
\begin{pmatrix}
\omega^{-1}&\omega^{-2}\\
0&\omega^{-1}
\end{pmatrix}.
\label{eq:ep-two-by-two-double-pole}
\end{equation}
No choice of eigenvector normalization removes the $\omega^{-2}$ term. The \term{Riesz moments} and \term{Laurent operators} introduced below generalize this coalescence of two poles, visible in the finite-dimensional model, to an isolated resonance cluster of a wave operator.

\section{Failure of individual residues}

Away from the exceptional point, write the Green operator associated with two simple poles as
\begin{equation}
\vG_{\mathrm{cl}}(\omega)
=
\frac{A_+}{\omega-\omega_+}
+\frac{A_-}{\omega-\omega_-}.
\label{eq:two-simple-poles}
\end{equation}
Here $A_+$ and $A_-$ are the \term{residue operators} of the individual poles. As the exceptional point is approached, \term{self-orthogonality} of the left and right eigenvectors can make $A_+$ and $A_-$ diverge separately. Their sum at fixed $\omega$ can nevertheless remain finite.

This divergence does not imply that the physical response must diverge. The coordinate components have become \term{ill conditioned} because two nearly parallel eigenvectors were chosen as coordinate axes. The object that survives at the exceptional point is not an individual mode label but the isolated two-state cluster as a whole.

\section{\term{Riesz contour moments}}

Let $\Gamma$ be a \term{counterclockwise} closed contour enclosing only the two poles and excluding all other poles and the continuous spectrum. Define two \term{contour moments} by
\begin{equation}
M_k
=
\frac{1}{2\pi\rmi}
\oint_\Gamma
\omega^k\vG(\omega)\,\rmd\omega,
\qquad
k=0,1.
\label{eq:contour_moments}
\end{equation}
When the simple poles are separated,
\begin{align}
M_0
&=
A_++A_-,
\label{eq:ep-moment-zero}
\\
M_1
&=
\omega_+A_+
+\omega_-A_-.
\label{eq:ep-moment-one}
\end{align}
Introduce the symmetric pole data
\begin{equation}
s_1=\omega_++\omega_- ,
\qquad
s_2=\omega_+\omega_-.
\end{equation}
Then Eq.~\eqref{eq:two-simple-poles} can be written as
\begin{equation}
\vG_{\mathrm{cl}}(\omega)
=
\frac{
M_1+(\omega-s_1)M_0
}{
\omega^2-s_1\omega+s_2
}.
\label{eq:pair-resolvent-moments}
\end{equation}
This representation is invariant under exchange of $+$ and $-$ and contains no divergence associated with the individual residues.

\section{Laurent operators}

At the exceptional point, where $\omega_+=\omega_-=\omegaEP$, the cluster Green operator has the expansion
\begin{equation}
\vG_{\mathrm{cl}}(\omega)
=
\frac{A_{-2}}{(\omega-\omegaEP)^2}
+\frac{A_{-1}}{\omega-\omegaEP}
+\vG_{\mathrm{reg}}(\omega).
\label{eq:ep-laurent}
\end{equation}
Equation~\eqref{eq:contour_moments} immediately gives
\begin{align}
A_{-1}
&=
M_0,
\label{eq:ep-a-minus-one}
\\
A_{-2}
&=
M_1-\omegaEP M_0.
\label{eq:ep-a-minus-two}
\end{align}
The operators $A_{-1}$ and $A_{-2}$ are defined without choosing a normalization for a \term{Jordan chain}.

Sandwich these operators between an observation operation and a source, and define
\begin{equation}
a_{-j}
=
\bra{O}A_{-j}\ket{S},
\qquad
j=1,2.
\end{equation}
The \term{inverse Fourier transform} of a \term{double pole} produces a time-domain response of the form
\begin{equation}
\vA_{\mathrm{cl}}(t)
=
\left(C_0+C_1t\right)
\rme^{-\rmi\omegaEP t}.
\label{eq:ep-time-response}
\end{equation}
The constants $C_0$ and $C_1$ are determined by $a_{-1}$, $a_{-2}$, and the transform convention. The term $t\rme^{-\rmi\omegaEP t}$ is not a fitting convenience but an unavoidable consequence of the double pole.

The relation among a double zero of the Jost function, the \term{Riesz--Laurent representation}, and the sourced response in black-hole scattering is given in Ref.~\cite{Morikawa:2026RieszLaurent}.

\begin{principlebox}
At an exceptional point, the representation in terms of individual modes becomes singular. The \term{contour-integral moments} of an isolated cluster, its Laurent operators, and its response for fixed source and observer remain finite and can be defined independently of mode labels and normalizations.
\end{principlebox}

\section{Geometric phase}

If the control parameters are varied adiabatically around a closed loop enclosing the exceptional point, the two states are exchanged after one circuit and return to their original labels after two. Rather than comparing only the signs of right eigenvectors, the \term{geometric phase} must be defined using the \term{biorthogonal connection} of left and right eigenvectors.

Resonance states cannot in general be normalized with the ordinary inner product, but complex scaling produces decaying states and allows a connection to be constructed over \term{parameter space}. The \term{geometric phase} and \term{topological invariants} associated with loops around exceptional points of quantum resonances are discussed in Ref.~\cite{Morikawa:2025inx}.

The \term{geometric phase} is not the same object as a local \term{Laurent coefficient}. The former is global information associated with a \term{parameter loop}; the latter is a local singularity in the frequency plane. Both, however, probe the geometry of a cluster rather than the normalization of an individual eigenvector.

\section{The real-frequency response can remain smooth}

Assume the following conditions near the exceptional point.

\begin{enumerate}[label=(\roman*),leftmargin=2.8em]
\item The two-pole cluster is isolated from the rest of the spectrum.
\item The resolvent on the \term{complementary subspace} is regular.
\item The poles do not reach the real frequencies being observed.
\end{enumerate}

Then Eq.~\eqref{eq:pair-resolvent-moments} is smooth at fixed real $\omega$, and transmission amplitudes and greybody factors can remain smooth as well. A divergence of individual residues or a \term{mode exchange} does not by itself imply an anomalously large scattering response. A separate limiting analysis is required if a pole reaches the real axis, the cluster collides with a continuous component, or the complementary subspace becomes singular.

\section{Conclusion of this chapter}

An exceptional point marks the limit of an individual-mode representation, not the limit of response theory. By moving to Riesz contour integrals and \term{Laurent operators}, one can define the double-pole time response, real-frequency scattering, and geometric phase as regular quantities in their appropriate senses.

\begin{researchproblem}[Higher-order exceptional points and continuous components]
Study an exceptional point of order three or higher, or the approach of a second-order exceptional point to a rotated continuous component. At minimum, determine the number of contour moments required, the degree of the \term{Laurent operators} and the associated time polynomial, and the conditions under which a contour can be chosen stably in a finite basis. Use the real-frequency response for a fixed source and observer, rather than the \term{condition number} of an individual eigenvector, as the final diagnostic of finiteness.
\end{researchproblem}


\part{Quantized Response}
\chapter{Quantization of the Master Field}
\label{chap:quantized-master-field}
\chapterepigraph{``What you have inherited from your fathers, earn it, in order to possess it.''}{Johann Wolfgang von Goethe, \textit{Faust I}, ``Night''}

\section{Reading the same action quantum mechanically}

Let $\Psi(t,x)$ be one of the master fields used in the classical theory, with action
\begin{equation}
I[\Psi]
=
\frac{1}{2}
\int \rmd t\rmd x
\left[
(\partial_t\Psi)^2
-(\partial_x\Psi)^2
-V(x)\Psi^2
\right].
\label{eq:quantum-master-action}
\end{equation}
Angular-momentum labels have been suppressed. The \term{canonical momentum} is
\begin{equation}
\Pi(t,x)
=
\partial_t\Psi(t,x).
\end{equation}
Under \term{quantization}, $\Psi$ and $\Pi$ become operators satisfying the \term{equal-time commutation relations}
\begin{align}
[\Psi(t,x),\Pi(t,x')]
&=
\rmi\delta(x-x'),
\label{eq:canonical-master-commutator}
\\
[\Psi(t,x),\Psi(t,x')]
&=
[\Pi(t,x),\Pi(t,x')]
=0.
\end{align}
We set $\hbar=1$. The \term{Heisenberg equation} is the same as the classical equation,
\begin{equation}
\left(
\partial_t^2-\partial_x^2+V
\right)\Psi
=0.
\label{eq:quantized-master-eom}
\end{equation}
For a linear field, the classical \term{propagator} therefore governs the quantum-field \term{commutator} as well.

\section{The retarded Green function is a commutator}

Let the quantum state be described by a \term{density operator} $\rho$. With the sign convention used in this book, define the retarded Green function by
\begin{equation}
G_R(t,x;t',x')
=
\rmi\theta(t-t')
\Tr\left(
\rho[\Psi(t,x),\Psi(t',x')]
\right).
\label{eq:quantum-retarded-green}
\end{equation}
For a linear equation, the \term{commutator} is a complex-valued \term{c-number} rather than a nontrivial operator, so the free-field $G_R$ is independent of the state $\rho$. The equal-time commutation relations fix the jump in the time derivative, and hence
\begin{equation}
\left(
\partial_t^2-\partial_x^2+V
\right)G_R(t,x;t',x')
=
\delta(t-t')\delta(x-x').
\label{eq:retarded-green-quantum-eom}
\end{equation}

If an external source $J$ is coupled through
\begin{equation}
H_{\mathrm{int}}(t)
=
-\int\rmd x\,J(t,x)\Psi(t,x),
\end{equation}
then the linear \term{Kubo response} is
\begin{equation}
\delta\langle\Psi(t,x)\rangle
=
\int\rmd t'\rmd x'\,
G_R(t,x;t',x')J(t',x').
\label{eq:kubo-master-response}
\end{equation}
This is the same solution obtained from the classical Green function in Chapter~\ref{chap:resolvent}.

\begin{principlebox}
The retarded Green function is the quantity that connects classical response to \term{quantum linear response}. For a linear field, its analytically continued QNM poles are the same as in the classical theory. Information about the quantum state enters not through the \term{commutator}, but through the anticommutator sector.
\end{principlebox}

\section{Two-point functions that contain the state}

To distinguish quantum states, use the \term{Wightman functions}
\begin{align}
G^>(1,2)
&=
\Tr\left(
\rho\Psi(1)\Psi(2)
\right),
\label{eq:wightman-greater}
\\
G^<(1,2)
&=
\Tr\left(
\rho\Psi(2)\Psi(1)
\right).
\label{eq:wightman-less}
\end{align}
Here $1=(t,x)$ and $2=(t',x')$ abbreviate spacetime points. Define the \term{spectral function} and \term{symmetric correlator} by
\begin{align}
\rho_\Psi(1,2)
&=
G^>(1,2)-G^<(1,2),
\label{eq:master-spectral-function}
\\
F_\Psi(1,2)
&=
\frac{1}{2}
\left[
G^>(1,2)+G^<(1,2)
\right].
\label{eq:master-symmetric-correlator}
\end{align}
Then $G_R=\rmi\theta(t-t')\rho_\Psi$. The function $\rho_\Psi$ measures what the system is capable of responding to, while $F_\Psi$ measures the fluctuations present in the state. In thermal equilibrium, they are related by the \term{fluctuation--dissipation relation} introduced in Chapter~\ref{chap:thermal-green}.

\section{Expansion in \term{scattering modes}}

Let $u_{\omega a}(t,x)$ be normalized real-frequency solutions. The label $a$ distinguishes independent \term{incoming channels}, such as modes incident from infinity or emerging from the past horizon. The field can be expanded as
\begin{equation}
\Psi(t,x)
=
\sum_a
\int_0^\infty\rmd\omega
\left[
a_{\omega a}u_{\omega a}(t,x)
+a^\dagger_{\omega a}u^*_{\omega a}(t,x)
\right].
\label{eq:master-mode-expansion}
\end{equation}
With \term{flux normalization},
\begin{equation}
[a_{\omega a},a^\dagger_{\omega' b}]
=
\delta_{ab}\delta(\omega-\omega').
\end{equation}
The reflection and transmission coefficients appear in the same classical scattering matrix; in the quantum theory they describe the transformation between channels of \term{creation and annihilation operators}.

QNMs must not be used as the fundamental operators in this expansion. A QNM is a complex-frequency resonance: a pole obtained by analytically continuing the complete real-frequency \term{scattering basis}. Preserving the \term{canonical commutation relations} of the quantum field requires the continuous components as well as the poles.

\section{Three vacuum states}

For a static black hole, the state depends on the time coordinate with respect to which positive frequency is defined. The vacuum associated with the static exterior time is the \term{Boulware state}; the state with no incoming particles from past infinity but a future \term{Hawking flux} is the \term{Unruh state}; and the state in which the two exteriors of an eternal black hole are in thermal equilibrium is the \term{Hartle--Hawking state}.

The essential point in this classification is to separate the equation of motion from the state. The same $G_R$ and scattering coefficients can be combined with different \term{occupation numbers}. Whether the state is regular at the horizon or carries a thermal flux at infinity is determined by the choice of state.

\begin{warningbox}
There is no unique ``\term{black-hole vacuum}.'' QNM frequencies are fixed by the background and radiation conditions, whereas particle numbers, noise, and Hawking flux depend on the quantum state and on the observer.
\end{warningbox}

\section{Conclusion of this chapter}

Canonical quantization of the master action makes the classical retarded Green function reappear as the field commutator. Poles and greybody factors are inherited from the classical theory, while the quantum state enters through the \term{Wightman functions} and \term{occupation numbers}. In the next chapter we show how analyticity near the horizon makes those occupation numbers thermal.

\chapter{Horizon Thermality}
\label{chap:horizon-thermality}
\chapterepigraph{``There are more things in heaven and earth, Horatio, than are dreamt of in your philosophy.''}{William Shakespeare, \textit{Hamlet}, Act I, Scene 5}
\chapterepigraph{``For now we see through a glass, darkly.''}{1 Corinthians 13:12, King James Version}

\section{The near-horizon geometry is Rindler spacetime}

For a \term{simple horizon} $r=r_h$ of a static spherically symmetric metric, define the \term{surface gravity} by
\begin{equation}
\kappa
=
\frac{1}{2}f'(r_h).
\label{eq:surface-gravity-static}
\end{equation}
Near the horizon, $f(r)\simeq2\kappa(r-r_h)$. Introducing the \term{proper distance} $\rho$ through
\begin{equation}
r-r_h
=
\frac{\kappa}{2}\rho^2,
\end{equation}
the $t$--$r$ part of the metric becomes
\begin{equation}
\rmd s^2
\simeq
-(\kappa\rho)^2\rmd t^2
+\rmd\rho^2.
\label{eq:near-horizon-rindler}
\end{equation}
Locally this is \term{Rindler spacetime}. The tortoise coordinate behaves as
\begin{equation}
\rstar
\simeq
\frac{1}{2\kappa}
\log(r-r_h),
\end{equation}
and the \term{outgoing null coordinate} $u=t-\rstar$ diverges at the horizon.

\section{\term{Kruskal coordinates} and analyticity}

Introduce a null coordinate regular on the future horizon,
\begin{equation}
U
=
-\rme^{-\kappa u}.
\label{eq:kruskal-U}
\end{equation}
In the exterior region, $U<0$. An outgoing \term{positive-frequency mode} with respect to the static time is
\begin{equation}
\rme^{-\rmi\omega u}
=
(-U)^{\rmi\omega/\kappa},
\label{eq:rindler-mode-U}
\end{equation}
which has a \term{branch point} at $U=0$. In the exterior $U<0$, take
$(-U)^{\rmi\omega/\kappa}=\exp[(\rmi\omega/\kappa)\operatorname{Log}(-U)]$.
Analytically continuing the \term{positive-frequency mode} through the lower half of the complex $U$ plane to $U>0$ gives
\begin{equation}
(-U)^{\rmi\omega/\kappa}
\longrightarrow
\rme^{-\pi\omega/\kappa}
U^{\rmi\omega/\kappa}.
\label{eq:horizon-branch-continuation}
\end{equation}
Combining this with the \term{negative-frequency component} continued through the upper half-plane and imposing \term{canonical normalization}, one obtains for the ratio of the \term{Bogoliubov coefficients} across the horizon
\begin{equation}
\left|
\frac{\beta_\omega}{\alpha_\omega}
\right|^2
=
\rme^{-2\pi\omega/\kappa}.
\label{eq:bogoliubov-thermal-ratio}
\end{equation}
Here $\alpha_\omega$ and $\beta_\omega$ are the Bogoliubov coefficients mixing positive and negative frequencies. Together with \term{canonical normalization}, this gives the bosonic \term{mean occupation number}
\begin{equation}
n_H(\omega)
=
\frac{1}{\rme^{\beta_H\omega}-1},
\qquad
\beta_H
=
\frac{2\pi}{\kappa}.
\label{eq:hawking-occupation}
\end{equation}
The \term{Hawking temperature} is therefore
\begin{equation}
T_H
=
\frac{\kappa}{2\pi}.
\label{eq:hawking-temperature}
\end{equation}
\cite{Hawking:1975vcx,BirrellDavies:1982}

\begin{principlebox}
The \term{thermal factor} arises from the near-horizon exponential map $U=-\rme^{-\kappa u}$ together with analyticity of positive-frequency modes. The exterior potential does not change the temperature; it changes the probability that the produced quanta reach infinity.
\end{principlebox}

\section{The greybody factor reappears}

Outgoing modes thermally populated near the horizon must transmit through the effective potential to reach infinity. Using the transmission probability $\Gamma_\ell(\omega)$ from Chapter~\ref{chap:resolvent}, the energy flux of a massless scalar field is
\begin{equation}
\frac{\rmd E}{\rmd t\rmd\omega}
=
\frac{1}{2\pi}
\sum_{\ell=0}^{\infty}
(2\ell+1)
\frac{
\omega\Gamma_\ell(\omega)
}{
\rme^{\beta_H\omega}-1
}.
\label{eq:hawking-greybody-flux}
\end{equation}
The \term{Planck distribution} of a planar blackbody is multiplied by scattering data from the black-hole exterior. Classical response theory is therefore used directly in computing quantum radiation.

For a black hole with angular momentum $a$ or charge $Q$, the frequency entering the \term{thermal factor} is the frequency measured with respect to the \term{horizon generator},
\begin{equation}
\widetilde\omega
=
\omega-m\Omega_H-q\Phi_H.
\label{eq:horizon-effective-frequency}
\end{equation}
Here $m$ is the \term{azimuthal quantum number}, $q$ is the field charge, $\Omega_H$ is the \term{horizon angular velocity}, and $\Phi_H$ is the \term{horizon electric potential}. In the regime $\widetilde\omega<0$, the flux sign of the transmission coefficient must be treated consistently with \term{superradiance}.

\section{The same temperature from \term{Euclidean time}}

Performing the \term{Wick rotation} $t=-\rmi\tau$ in Eq.~\eqref{eq:near-horizon-rindler} gives
\begin{equation}
\rmd s_E^2
\simeq
\rmd\rho^2
+\kappa^2\rho^2\rmd\tau^2.
\end{equation}
Avoiding a \term{conical singularity} at $\rho=0$ requires
\begin{equation}
\tau
\sim
\tau+\frac{2\pi}{\kappa}.
\label{eq:euclidean-period}
\end{equation}
The period of \term{imaginary time} is therefore $\beta_H$. The calculation using \term{Bogoliubov coefficients} and the calculation using \term{Euclidean regularity} are distinct, but both read the same near-horizon geometry.

\section{Area entropy and the first law}

In this section we set $G=c=\hbar=k_B=1$. If the event-horizon area is $A_H$, the \term{Bekenstein--Hawking entropy} is
\begin{equation}
S_{\mathrm{BH}}
=
\frac{A_H}{4}.
\label{eq:bekenstein-hawking-entropy}
\end{equation}
\cite{Bekenstein:1973ur,Hawking:1975vcx} Restoring conventional units gives $S_{\mathrm{BH}}=k_Bc^3A_H/(4G\hbar)$. The \term{first law} relating nearby stationary \term{Kerr--Newman black holes} is
\begin{equation}
\delta M
=
\frac{\kappa}{8\pi}\delta A_H
+\Omega_H\delta J
+\Phi_H\delta Q
=
T_H\delta S_{\mathrm{BH}}
+\Omega_H\delta J
+\Phi_H\delta Q.
\label{eq:black-hole-first-law}
\end{equation}
\cite{Bardeen:1973gs} This is not a derivation of the microscopic number of states. What response theory needs is the conversion between fluxes of energy, angular momentum, and charge entering the horizon and the resulting changes in the stationary solution and its entropy.

\paragraph{A signpost toward the Ryu--Takayanagi formula.}
When classical gravity in a static \term{asymptotically AdS} spacetime describes a \term{boundary quantum theory}, the \term{entanglement entropy} of a spatial boundary region $A$ is
\begin{equation}
S_A
=
\frac{\operatorname{Area}(\gamma_A)}{4G_N}.
\label{eq:ryu-takayanagi-signpost}
\end{equation}
Here $\gamma_A$ is a \term{codimension-two minimal surface} in the \term{bulk} with the same boundary as $A$ and \term{homologous} to $A$ \cite{Ryu:2006bv}. This extends Eq.~\eqref{eq:bekenstein-hawking-entropy} in the sense that an area is interpreted as an entropy. Because this book does not construct the \term{AdS/CFT} dictionary, however, the formula is not used in later derivations. Time-dependent versions, quantum corrections, and the \term{entanglement wedge} are deferred to research problems only when a concrete response question requires them.

\section{Differences among quantum states}

The Boulware state contains no particles for static observers at infinity but is singular at the horizon. The Hartle--Hawking state is regular on both future and past horizons and places the exterior in a \term{thermal bath} at temperature $T_H$. The Unruh state corresponds to a black hole formed by collapse: incoming modes from past infinity are in their vacuum state while a Hawking flux emerges at future infinity.

Equation~\eqref{eq:hawking-temperature} gives the local horizon temperature scale, but it does not imply that every two-point function is in the same global thermal-equilibrium state. In particular, the Unruh state is not a thermal equilibrium state of the entire static exterior.

\begin{warningbox}
The \term{Hawking temperature}, the particle flux at infinity, and the response of a \term{local detector} are not the same observable. The temperature is fixed by the \term{surface gravity}; the flux additionally depends on the \term{greybody factor} and the quantum state; and a \term{detection rate} also depends on the detector trajectory and coupling.
\end{warningbox}

\section{Conclusion of this chapter}

The exponential relation between horizon-regular and static coordinates mixes positive frequencies and gives the temperature $T_H=\kappa/(2\pi)$. Classical scattering in the exterior filters the resulting thermal occupation numbers through the \term{greybody factor}. In the next chapter we formulate thermality not in terms of particle numbers but through the KMS condition for two-point functions.

\chapter{Thermal Green Functions}
\label{chap:thermal-green}
\chapterepigraph{``The energy of the universe is constant. The entropy of the universe strives toward a maximum.''}{Rudolf Clausius, ``Concerning Several Conveniently Applicable Forms for the Main Equations of the Mechanical Heat Theory'' (1865)}

\section{The KMS condition}

Let $H$ be the Hamiltonian generating time evolution, and let the \term{equilibrium state} at \term{inverse temperature} $\beta$ be
\begin{equation}
\rho_\beta
=
\frac{\rme^{-\beta H}}{Z},
\qquad
Z
=
\Tr\rme^{-\beta H}.
\label{eq:gibbs-state}
\end{equation}
For a \term{Hermitian operator} $O(t)$, define
\begin{align}
G^>(t)
&=
\Tr\left(
\rho_\beta O(t)O(0)
\right),
\\
G^<(t)
&=
\Tr\left(
\rho_\beta O(0)O(t)
\right).
\end{align}
Cyclicity of the \term{trace}, together with \term{imaginary-time evolution}, gives
\begin{equation}
G^>(t)
=
G^<(t+\rmi\beta).
\label{eq:kms-time}
\end{equation}
This is the \term{Kubo--Martin--Schwinger (KMS) condition}. Even in field theory, where a finite-dimensional Gibbs matrix cannot be written, analyticity of \term{correlation functions} together with Eq.~\eqref{eq:kms-time} can be taken as the definition of \term{thermal equilibrium}.

Using the Fourier transform
\begin{equation}
G(\omega)
=
\int_{-\infty}^{\infty}\rmd t\,
\rme^{+\rmi\omega t}G(t),
\label{eq:thermal-fourier}
\end{equation}
the \term{KMS condition} becomes the \term{detailed-balance} relation
\begin{equation}
G^>(\omega)
=
\rme^{\beta\omega}G^<(\omega).
\label{eq:kms-frequency}
\end{equation}

\section{\term{Lehmann representation}}

Let $H\ket{n}=E_n\ket{n}$ and $O_{mn}=\bra{m}O\ket{n}$. Inserting a \term{complete set} into the two-point functions gives
\begin{align}
G^>(\omega)
&=
\frac{2\pi}{Z}
\sum_{m,n}
\rme^{-\beta E_n}
|O_{mn}|^2
\delta\bigl(\omega-E_m+E_n\bigr),
\label{eq:lehmann-greater}
\\
G^<(\omega)
&=
\frac{2\pi}{Z}
\sum_{m,n}
\rme^{-\beta E_m}
|O_{mn}|^2
\delta\bigl(\omega-E_m+E_n\bigr).
\label{eq:lehmann-less}
\end{align}
On the support of the delta function, $E_m-E_n=\omega$, so each term separately satisfies $G^>(\omega)=\rme^{\beta\omega}G^<(\omega)$. At \term{finite volume}, the spectrum is a sum of \term{discrete lines}; in the \term{thermodynamic limit}, or after \term{coarse graining} at finite resolution, it becomes a smooth \term{spectral function}.

This representation shows that a thermal Green function is not merely an abstract analytic function: it is a sum over matrix elements of observables weighted by the \term{density of states}. The ETH of the next chapter assumes a smooth envelope for the individual $O_{mn}$ appearing in Eq.~\eqref{eq:lehmann-greater} and asks under what conditions the same \term{detailed balance} emerges from a single high-energy eigenstate.

\section{Spectral function and fluctuations}

Define the \term{spectral function} and \term{symmetric correlator} by
\begin{align}
\rho_O(\omega)
&=
G^>(\omega)-G^<(\omega),
\label{eq:thermal-spectral-density}
\\
F_O(\omega)
&=
\frac{1}{2}
\left[
G^>(\omega)+G^<(\omega)
\right].
\label{eq:thermal-symmetric-density}
\end{align}
Equation~\eqref{eq:kms-frequency} then gives
\begin{equation}
F_O(\omega)
=
\frac{1}{2}
\coth\left(\frac{\beta\omega}{2}\right)
\rho_O(\omega).
\label{eq:fdt-rho}
\end{equation}

The retarded Green function used in this book is
\begin{equation}
G_R(t)
=
\rmi\theta(t)
\Tr\left(
\rho_\beta[O(t),O(0)]
\right).
\end{equation}
For its real-frequency boundary value,
\begin{equation}
\rho_O(\omega)
=
2\im G_R(\omega),
\label{eq:spectral-im-retarded}
\end{equation}
so the \term{fluctuation--dissipation relation} is
\begin{equation}
F_O(\omega)
=
\coth\left(\frac{\beta\omega}{2}\right)
\im G_R(\omega).
\label{eq:fluctuation-dissipation}
\end{equation}
The \term{dissipative part} of the response determines the quantum and thermal fluctuations in equilibrium.

\begin{warningbox}
The sign in Eq.~\eqref{eq:spectral-im-retarded} depends on the conventions for the retarded function and the Fourier transform. When importing a formula from the literature, check as one package the factor of $\rmi$ in front of $G_R$, the sign in $\rme^{\pm\rmi\omega t}$, and the ordering used in the \term{spectral function}.
\end{warningbox}

\section{Absorption and emission}

Couple a weak periodic external source $j(t)$ through $-j(t)O(t)$. For $\omega>0$, the absorbed power per unit time is proportional to $\rho_O(\omega)$, equivalently to $\im G_R(\omega)$. By contrast, \term{spontaneous and stimulated emission} from the state is measured by $G^<(\omega)$. The \term{KMS condition} requires
\begin{equation}
\frac{G^<(\omega)}{G^>(\omega)}
=
\rme^{-\beta\omega}.
\label{eq:detailed-balance-ratio}
\end{equation}

For a black hole, the classical \term{absorption cross section} or greybody factor gives the real-axis dissipation contained in $\rho_O$. In the Hartle--Hawking state, the \term{KMS condition} relates absorption to emission. In the Unruh state, the outgoing sector carries the \term{Hawking occupation number}, but the full state is not in global \term{thermal equilibrium}.

\begin{principlebox}
Classical absorption, quantum emission, and \term{equilibrium noise} are not three unrelated tables of data. The imaginary part of the retarded Green function and the KMS condition organize them into the same spectral function.
\end{principlebox}

\section{Do not confuse QNMs with thermal poles}

Finite-temperature calculations involve the \term{Euclidean frequencies}
\begin{equation}
\omega_n^{E}
=
\frac{2\pi n}{\beta},
\qquad
n\in\bbZ.
\label{eq:bosonic-matsubara}
\end{equation}
These are \term{Matsubara frequencies} arising from periodicity in imaginary time; they are not QNM frequencies determined by radiation boundary conditions. Appropriate analytic continuation of the \term{Euclidean correlator} gives the retarded function, and QNMs then appear as poles on a different sheet of that real-time object.

Thus, even if a sequence of \term{Matsubara points} resembles a sequence of QNM poles in a plot, the two must not be identified. The former are sampling points associated with a thermal state; the latter are resonances of real-time response. Relating them requires uniqueness of the analytic continuation together with control of the high-frequency behavior.

\section{Thermal equilibrium in a rotating system}

For a stationary rotating system, the density operator has the schematic form
\begin{equation}
\rho
\mathrel{\propto}
\exp\left[
-\beta
\left(
H-\Omega_H J-\Phi_H Q
\right)
\right].
\label{eq:rotating-gibbs}
\end{equation}
The KMS factor depends on the effective frequency $\widetilde\omega$ defined in Eq.~\eqref{eq:horizon-effective-frequency}. In asymptotically flat Kerr spacetime, \term{superradiant modes} obstruct the construction of an ordinary Hartle--Hawking state covering the entire exterior. One must distinguish the local use of a thermal factor from the existence of a global \term{Gibbs state}.

\section{Conclusion of this chapter}

A thermal state is characterized by the KMS condition, and the \term{fluctuation--dissipation relation} connects state-dependent fluctuations to the state-independent retarded response. Classical black-hole absorption and quantum emission lie on opposite sides of this relation. In the next chapter we introduce ETH, which attempts to explain thermal equilibrium in terms of matrix elements within a single energy eigenstate.

\chapter{The Eigenstate Thermalization\\ Hypothesis}
\label{chap:eth}
\chapterepigraph{%
  \jp{井蛙不可以語於海者，拘於虛也；\\
  夏蟲不可以語於冰者，篤於時也；\\
  曲士不可以語於道者，束於教也。}%
}{\textit{Zhuangzi}, Outer Chapters, ``Autumn Floods''}

\section{Can a single eigenstate be thermal?}

The \term{KMS condition} in the previous chapter characterizes correlation functions in thermal equilibrium. An \term{isolated quantum system}, however, evolves \term{unitarily} and remains in a \term{pure state}. The \term{eigenstate thermalization hypothesis} (\term{ETH}) seeks the reason why observables probing only a small number of degrees of freedom can nevertheless approach thermal-equilibrium values in the Hamiltonian eigenstates themselves \cite{Deutsch:1991,Srednicki:1994,Rigol:2008,DAlessio:2016}.
Let
\begin{equation}
H\ket{n}
=
E_n\ket{n}
\end{equation}
and denote matrix elements of an observable $O$ by $O_{mn}=\bra{m}O\ket{n}$. The standard \term{ETH ansatz} is
\begin{equation}
O_{mn}
=
\overline O(\overline E)\delta_{mn}
+\rme^{-S(\overline E)/2}
f_O(\overline E,\omega)R_{mn},
\label{eq:eth-ansatz}
\end{equation}
where
\begin{equation}
\overline E
=
\frac{E_m+E_n}{2},
\qquad
\omega
=
E_m-E_n.
\end{equation}
Here $S(E)$ is the \term{microcanonical entropy}; $\overline O$ and $f_O$ are smooth functions on macroscopic energy scales; and $R_{mn}$ is an irregular quantity with zero mean and unit variance.

The diagonal term states that the static expectation value in a single eigenstate approaches the \term{microcanonical average}. The factor $\rme^{-S/2}$ in the off-diagonal term makes each individual transition small, but when combined with an exponentially large number of final states it produces a finite dynamical correlation function. \term{ETH} is not a theorem valid for every Hamiltonian and every observable. \term{Integrability}, \term{localization}, \term{quantum many-body scars}, and symmetries all modify its domain of validity.

\section{From matrix elements to the spectral function}

The Wightman function in a single \term{energy eigenstate} $\ket{n}$ is
\begin{equation}
G_n^>(t)
=
\bra{n}O(t)O(0)\ket{n}
=
\sum_m
\rme^{-\rmi(E_m-E_n)t}
|O_{nm}|^2.
\label{eq:eigenstate-wightman}
\end{equation}
Fourier transformation gives
\begin{equation}
G_n^>(\omega)
=
2\pi
\sum_m
|O_{nm}|^2
\delta\left(
\omega-E_m+E_n
\right).
\label{eq:eigenstate-spectral-sum}
\end{equation}
Insert Eq.~\eqref{eq:eth-ansatz} and approximate the \term{density of states} smoothly near $E_n$. Using the thermodynamic relation
\begin{equation}
\beta(E)
=
\frac{\partial S(E)}{\partial E},
\label{eq:microcanonical-beta}
\end{equation}
and expanding the entropy for small $\omega$, the ratio of positive- and negative-frequency components becomes
\begin{equation}
\frac{G_n^>(\omega)}{G_n^<(\omega)}
\simeq
\rme^{\beta(E_n)\omega}.
\label{eq:eth-detailed-balance}
\end{equation}
Thus a \term{coarse-grained} two-point function in a single high-energy eigenstate can reproduce the KMS \term{detailed-balance} relation of the previous chapter.

\begin{principlebox}
The place where \term{ETH} enters response theory is not the QNM frequency itself but the matrix elements of the observable. The \term{diagonal elements} determine static thermodynamic behavior, while the envelope $f_O$ of the \term{off-diagonal elements} determines the spectral function and dissipation.
\end{principlebox}

\section{Resolve symmetries first}

If a \term{conserved charge} $Q$ commutes with the Hamiltonian, the \term{Hilbert space} decomposes into \term{sectors of fixed charge}. Standard ETH must be formulated within a single \term{irreducible sector}, without mixing different charge sectors. If the observable itself carries charge, \term{selection rules} leave only matrix elements between appropriate sectors. Thus the $R_{mn}$ in Eq.~\eqref{eq:eth-ansatz} must not be treated as independent random variables without regard to symmetry.

When several \term{conserved charges} fail to commute, giving a \term{non-Abelian symmetry}, there is no basis that diagonalizes all charges simultaneously, and degeneracies associated with \term{symmetry multiplets} arise. \term{Non-Abelian ETH} incorporates an \term{approximately microcanonical subspace} and the \term{representation structure} of the symmetry into the ansatz in order to compare time averages and thermal averages of \term{local observables} \cite{Murthy:2023naETH}. We do not need the details here; we use only the fact that even the standard prescription ``fix the symmetry sector'' itself requires modification.

\section{Higher-form symmetries change the class of observables}

An ordinary \term{global symmetry} acts on pointlike \term{charged operators}, whereas a $p$-form symmetry acts on charged operators extended over $p$ dimensions. In this situation, thermalization of \term{local operators} does not guarantee that line- or surface-supported operators thermalize to the same ensemble.

Fukushima and Hamazaki showed, under reasonable assumptions, that a $p$-form symmetry in a $(d+1)$-dimensional quantum field theory can violate ETH for many nontrivial $(d-p)$-dimensional observables \cite{Fukushima:2023svf}. In a $(2+1)$-dimensional $\bbZ_2$ \term{lattice gauge theory}, local \term{plaquette operators} thermalize, while \term{extended operators} that excite a \term{magnetic dipole} relax to a \term{generalized Gibbs ensemble} containing the $\bbZ_2$ \term{one-form charge}.

The lesson is that one should not ask a single yes-or-no question, ``Does the system satisfy ETH?'' One must state which symmetries are fixed, what dimensionality and charge the measured observable carries, and which statistical ensemble provides the comparison. The same distinction is required in black-hole systems with \term{gauge fields} and \term{boundary degrees of freedom}; local ETH should not be extrapolated to extended observables without this analysis.

\begin{warningbox}
Symmetry-induced modifications of ETH do not by themselves establish thermalization or nonthermalization in a particular \term{microscopic black-hole model}. The conserved charges, support of the observable, boundary conditions, and comparison ensemble must first be identified explicitly.
\end{warningbox}

\section{Where quantum field theory meets black-hole response}

In field theory, ETH for \term{local operators} is not the same statement as ETH for the state of a \term{finite region}. In a \term{conformal field theory}, it has been shown that if \term{local probes} in high-energy \term{primary eigenstates} satisfy ETH, the statement can be extended to observables supported on a finite region and measuring finite energy \cite{Lashkari:2018CFTETH}. This is an example of the additional assumptions needed to move from local matrix elements to the response of a \term{subsystem}.

Black-hole linear response is described by the retarded Green function $G_R$, whose imaginary part on the real axis is the spectral function of the operator $O$. If a \term{microscopic Hamiltonian} and the operator coupled to the external field satisfy ETH in the appropriate symmetry sector, the following correspondences are expected.

\begin{enumerate}[label=(\roman*),leftmargin=2.8em]
\item The derivative of the entropy $S(E)$ determines an \term{effective temperature}.
\item $f_O(E,\omega)$ determines the smooth envelope of absorption and emission.
\item \term{Coarse graining} the discrete lines at \term{finite entropy} produces a smooth spectral function.
\item Poles of the analytic continuation of the corresponding retarded function can approach the classical response poles.
\end{enumerate}

This is not a derivation of ETH from classical QNMs. A QNM is a pole in the analytic continuation of a \term{coarse-grained} retarded response, whereas ETH is a hypothesis about microscopic matrix elements. Connecting them requires separate control of symmetry-preserving coarse graining, the operator correspondence, and \term{finite-entropy} errors.

\section{The physics-minimum scope}

The \term{random-matrix models} of ETH, \term{quantum many-body scars}, integrable systems, \term{many-body localization}, and \term{out-of-time-order correlators} are each independent subjects. The main line of this book requires only the ansatz \eqref{eq:eth-ansatz}, the spectral sum \eqref{eq:eigenstate-spectral-sum}, its connection to the KMS condition, and the modifications of the domain of validity caused by symmetries and the choice of observables.

Within this scope, ETH is not a survey of the information problem but the endpoint that connects retarded response to microscopic matrix elements. A new symmetry or microscopic model should return to the main text only when it actually changes the response function or the statistical ensemble against which it is compared.

\section{Conclusion of this chapter}

ETH is a hypothesis about matrix elements that allows thermal two-point functions to emerge from a single high-energy eigenstate. Its meeting point with black-hole response is not the numerical value of a QNM frequency but the spectral function, detailed balance, and symmetry-preserving coarse graining. Commuting charges, noncommuting charges, and \term{higher-form symmetries} change the sectors, observables, and statistical ensembles relevant to any claim of thermalization.

\begin{researchproblem}[Symmetry-resolved ETH response]
Choose a finite-entropy gauge theory or \term{holographic Hamiltonian}. Identify its ordinary global charges, higher-form charges, and the transformation properties of the local and \term{extended operators} to be measured. Measure matrix elements sector by sector, and distinguish observables for which standard ETH holds from those requiring a \term{generalized Gibbs ensemble}. Then compare the real-axis \term{spectral density} of the coarse-grained retarded Green function with its analytically continued poles. Do not assume agreement with a classical black hole in advance; report as results the frequency window, \term{coarse-graining width}, and entropy dependence of the errors for which the comparison holds.
\end{researchproblem}

\chapter{Response Invariants}
\label{chap:response-invariants}
\chapterepigraph{``In discussion it is not so much weight of authority as force of argument that should be demanded.''}{Marcus Tullius Cicero, \textit{De Natura Deorum}, I.10, trans.\ H.\ Rackham}

\section{What survives a change of representation}

The choice of master field, coordinate, basis, \term{complex-scaling angle}, and eigenvector \term{normalization} is necessary for a calculation, but none of these choices is by itself an observable. We conclude the book by organizing the hierarchy of quantities that make up a response.

Write the frequency-domain equation as
\begin{equation}
\vL(\omega)\ket{\Psi}
=
\ket{S}.
\end{equation}
Under an \term{invertible field transformation} $B(\omega)$ and an invertible recombination $A(\omega)$ of the equations, let
\begin{align}
\ket{\Psi'}
&=
B\ket{\Psi},
\\
\vL'
&=
A\vL B^{-1},
\\
\ket{S'}
&=
A\ket{S}.
\end{align}
Then
\begin{equation}
\vG'
=
B\vG A^{-1}.
\label{eq:green-field-redefinition}
\end{equation}
If the observation operation is transformed as $\bra{O'}=\bra{O}B^{-1}$, then
\begin{equation}
\bra{O'}\vG'\ket{S'}
=
\bra{O}\vG\ket{S}.
\label{eq:physical-amplitude-invariant}
\end{equation}
Thus, although components of the \term{Green operator} change, the amplitude connecting the same physical source to the same observer does not.

\section{What should be compared}

When comparing results from different calculations, do not mix the following levels.

\paragraph{Representation.}
To compare $\rstar$, $V$, $\psi_n$, or basis coefficients, match the coordinate, field definition, and \term{normalization}.

\paragraph{Spectrum.}
To compare QNM poles, branch cuts, or a \term{Riesz cluster}, match the boundary conditions, \term{Riemann surface}, and \term{contour}.

\paragraph{Classical response.}
To compare $\bra{O}\vG\ket{S}$, the scattering matrix, or a greybody factor, match the physical source, observation operation, and flux normalization.

\paragraph{Quantum state.}
To compare $G^>$, $G^<$, particle numbers, or noise, match the state, time-evolution generator, and observer.

\paragraph{Thermalization.}
To compare $S(E)$, $f_O(E,\omega)$, or \term{charge sectors}, match the \term{coarse graining}, symmetries, and statistical ensemble.

QNM frequencies are stable under \term{admissible} field redefinitions, although a pole crossing a \term{branch cut} requires a specification of the \term{sheet}. An \term{operator-valued residue} transforms with the representation, whereas the matrix element in Eq.~\eqref{eq:physical-amplitude-invariant} is invariant. At an exceptional point, do not compare individual residues; compare $M_0$, $M_1$, and sourced \term{Laurent coefficients} defined with the same contour.

\section{Dispersion relations from causality}

A retarded Green function vanishes for $t<0$. This \term{causality} makes $G_R(\omega)$ analytic in the \term{upper half-plane}. For matrix elements that decay sufficiently rapidly at infinity, the real and imaginary parts on the real axis satisfy
\begin{align}
\re G_R(\omega)
&=
\frac{1}{\pi}
\mathop{\mathrm{P}}
\int_{-\infty}^{\infty}
\frac{\im G_R(\omega')}{\omega'-\omega}
\,\rmd\omega',
\label{eq:kramers-kronig-real}
\\
\im G_R(\omega)
&=
-\frac{1}{\pi}
\mathop{\mathrm{P}}
\int_{-\infty}^{\infty}
\frac{\re G_R(\omega')}{\omega'-\omega}
\,\rmd\omega'.
\label{eq:kramers-kronig-imag}
\end{align}
Here $\mathrm P$ denotes the \term{Cauchy principal value}. If constants or polynomial terms remain at high frequency, appropriate \term{subtractions} are required.

These \term{dispersion relations} express, directly on the real axis, why poles alone are insufficient and continuous components are also needed. One cannot arbitrarily change dispersion while leaving absorption unchanged. The relations also provide a numerical consistency check comparing the real and imaginary parts of a computed Green function.

\section{One equation from classical response to quantum statistics}

The main line of this book can be summarized by the following \term{causal amplitude}:
\begin{equation}
\vA_R(\omega)
=
\bra{O}
\left(
H-(\omega+\rmi0)^2
\right)^{-1}
\ket{S}.
\label{eq:book-central-amplitude}
\end{equation}
The geometry enters through $H$; the horizon through the $+\rmi0$ boundary value and the radiation conditions; and the external forcing and detector through $\ket S$ and $\bra O$. Analytic continuation reveals QNMs and exceptional points, while the imaginary part on the real axis gives absorption. After quantization, the same object becomes a commutator; specifying a state adds the KMS condition and fluctuation--dissipation relation. Returning to microscopic eigenstates, its \term{spectral density} can be represented in terms of ETH matrix elements.

\begin{principlebox}
Preserve the causal response, not a particular representation. The tortoise coordinate, Jost solutions, QNMs, complex scaling, Riesz clusters, KMS condition, and ETH are tools for reading Eq.~\eqref{eq:book-central-amplitude} in different regimes.
\end{principlebox}

\newpage
\section{The boundary of the physics minimum}

The connection from nonlinear merger to ringdown, the \term{self-force}, connections to \term{numerical relativity}, the \term{AdS/CFT} dictionary, \term{microscopic states} of \term{quantum gravity}, and \term{information recovery} are all important subjects. But adding one survey chapter on each would not deepen the causal response developed here. They should enter a future edition only when they can be reduced to a concrete source, observable, Green function, or matrix element.

For the same reason, lists of numerical methods and tables of known QNM frequencies have been kept out of the main text. Necessary calculations are placed in research problems, where the emphasis is on diagnostic criteria rather than on quoted results. When that research is completed, only results that alter the main line should return to the text; the rest should remain as updated research problems.

\section{Conclusion}

A black hole is a scatterer that absorbs classical fields, an open system with resonances, and a horizon that produces thermal fluctuations in quantum fields. These need not be treated as unrelated metaphors. Once the causal Green function is placed at the center, classical-field response, the \term{non-self-adjoint spectrum}, \term{quantum linear response}, and thermalization become successive descriptions of the same problem.

The physics minimum adopted in this book is not merely a restrictive device for keeping the text short. It is a choice designed to keep sight of which calculations actually reach observables.

\chapterepigraph{``And thence we came forth to see again the stars.''}{Dante Alighieri, \textit{Inferno}, Canto XXXIV, line 139}


\appendix
\chapter{Notation and Conventions}
\label{app:conventions}
\chapterepigraph{``The last thing we decide on in writing a book is what shall be the first we put in it.''}{Attributed to Blaise Pascal, \textit{Pens\'ees}, Brunschvicg 19, trans.\ C.\ Kegan Paul}

This appendix collects the conventions used throughout the book. When comparing formulas from different references, one should first check, in order, the \term{time convention} for frequency, the ordering in the \term{Wronskian}, and the sign convention for the \term{resolvent}.

\section{Spacetime and units}

The metric signature is $(-,+,+,+)$, and we use \term{geometrized units} $G=c=1$. A \term{static spherically symmetric metric} is written as
\begin{equation}
\rmd s^2
=
-f(r)\,\rmd t^2
+f(r)^{-1}\rmd r^2
+r^2\rmd\Omega^2.
\end{equation}
The tortoise coordinate is defined by $\rmd\rstar/\rmd r=f(r)^{-1}$.

\section{\term{Fourier convention}}

The time dependence is
\begin{equation}
\Psi(t,x)
=
\int
\frac{\rmd\omega}{2\pi}
\,
\rme^{-\rmi\omega t}
\psi(\omega,x).
\end{equation}
With this convention, $\im\omega<0$ represents temporal decay. A right-moving wave is $\rme^{-\rmi\omega(t-x)}$, while a left-moving wave is $\rme^{-\rmi\omega(t+x)}$.

\section{Operators and resolvents}

The \term{spatial operator}, \term{frequency operator pencil}, and \term{resolvents} are defined by
\begin{align}
H
&=
-\frac{\rmd^2}{\rmd x^2}+V(x),
\\
\vL(\omega)
&=
H-\omega^2,
\\
\vR(z)
&=
(H-z)^{-1},
\\
\vG(\omega)
&=
\vL(\omega)^{-1}.
\end{align}
Some references define the resolvent as $(z-H)^{-1}$, which reverses the overall sign of the \term{Riesz projection}. To avoid this ambiguity, the \term{contour moments} in this book are defined directly by Eq.~\eqref{eq:contour_moments}.

\section{Jost solutions and the Wronskian}

The left and right \term{Jost solutions} are normalized by
\begin{align}
f_-(x,\omega)
&\sim
\rme^{-\rmi\omega x},
&&x\longrightarrow-\infty,
\\
f_+(x,\omega)
&\sim
\rme^{+\rmi\omega x},
&&x\longrightarrow+\infty.
\end{align}
The \term{Wronskian} is
\begin{equation}
\vW(\omega)
=
f_-\frac{\partial f_+}{\partial x}
-
\frac{\partial f_-}{\partial x}f_+.
\end{equation}
With this ordering, the \term{Green kernel} carries the minus sign shown in Eq.~\eqref{eq:green_jost_formula}.

\section{Abbreviations}

In the main text, \term{quasinormal mode} is abbreviated QNM, \term{complex scaling method} as CSM, and \term{exceptional point} as EP. Regge--Wheeler and Zerilli are occasionally abbreviated RW and Z, respectively.

\section{Quantum two-point functions}

The retarded Green function is defined by
\begin{equation}
G_R(t)
=
\rmi\theta(t)\langle[O(t),O(0)]\rangle.
\end{equation}
Its Fourier transform uses $\int\rmd t\,\rme^{+\rmi\omega t}$. With this convention, the \term{spectral function} $\rho_O=G^>-G^<$ satisfies $\rho_O=2\im G_R$. We set $\hbar=1$.

\section{Cautions in complex analysis}

Because $z=\omega^2$, the \term{energy plane} and \term{frequency plane} form a double cover. For square roots, logarithms, and the \term{analytic continuation} of \term{Jost solutions}, one must fix the \term{branch cut} and the \term{sheet}. In numerical calculations it is desirable to record not only the point in the complex plane but also the path of analytic continuation used to reach it.

When encircling an exceptional point, one must specify the path in control-parameter space, the \term{parallel-transport convention} for eigenvectors, and the \term{biorthogonal inner product}. \term{Eigenvalue permutation}, the sign of an eigenvector, and the \term{geometric phase} are related but are not identical statements.

\chapter*{About This Revised Draft}
\addcontentsline{toc}{chapter}{About This Revised Draft}

This revised edition follows the policy that ``research to be done next belongs in research problems, and material unnecessary to the main line should be removed.'' The main text takes a single path: construct the classical Green function from the tortoise coordinate, read its poles and continuous components, connect the same retarded function to quantum commutators, the KMS condition, and ETH matrix elements, and finally return to response invariants. The intended readership is stated explicitly as senior undergraduates. As a physics minimum, we have added coordinates regular at the future horizon, the non-radiative sectors of gravitational perturbations, one complete chain from a source to a waveform and flux, Kerr radiation conditions and \term{superradiance}, and the \term{area entropy} and \term{first law}. The \term{Ryu--Takayanagi formula} is not used in the main derivation; it appears only as a brief signpost from the area law toward boundary theory.

There are no conventional end-of-chapter exercises. Instead, we pose research problems on extremal horizons, long-range Jost functions, numerical implementation of Kerr QNMs, a uniform approximation joining poles and tails, continuous response in de~Sitter space, higher-order exceptional points, and symmetry-resolved ETH. These are not checks of comprehension but starting points for moving beyond the assumptions of the text.

Future revisions will not aim to add chapters for their own sake. A new result will return to the main text only if it changes the causal response connecting source and observer or separates a new invariant from representation-dependent quantities. Comparisons of numerical methods, tables of known frequencies, and broad surveys will remain in research problems or a separate implementation note unless they alter the main line. Roughly one hundred pages of main text excluding references remains an editorial target, not a strict bound that would justify disrupting the order of definitions.


\backmatter
\bibliographystyle{utphys}
\bibliography{ref}
\printindex

\end{document}
